\documentclass[11pt]{article}
\usepackage[a4paper,margin=1in]{geometry}
\usepackage[T1]{fontenc}
\usepackage{amsmath,amssymb,amsthm}
\usepackage{array,booktabs,tabularx}
\usepackage{enumitem}
\usepackage{pgfplots}
\pgfplotsset{compat=1.17}
\usepackage{xcolor}
\definecolor{linkblue}{RGB}{24,69,102}
\usepackage[colorlinks=true,allcolors=linkblue]{hyperref}

\setlist{itemsep=2pt,topsep=4pt}
\setlength{\emergencystretch}{2em}

\newcommand{\dd}{\mathrm{d}}
\newcommand{\vt}{v_t}
\newcommand{\Kn}{\mathrm{Kn}}
\newcommand{\sgn}{\operatorname{sgn}}
\newcommand{\avg}[1]{\langle #1 \rangle}
\newcommand{\bu}{\mathbf{u}}
\newcommand{\bv}{\mathbf{v}}
\newcommand{\bw}{\mathbf{w}}
\newcommand{\bc}{\mathbf{c}}
\newcommand{\bq}{\mathbf{q}}
\newcommand{\bx}{\mathbf{x}}
\newcommand{\bE}{\mathbf{E}}
\newcommand{\bB}{\mathbf{B}}
\newcommand{\bF}{\mathbf{F}}
\newcommand{\bb}{\mathbf{b}}
\newcommand{\Hil}{\mathcal{H}}
\newcommand{\PV}{\mathrm{P.V.}\!\!}
\newcommand{\He}{\mathrm{He}}
\newcommand{\kF}{\kappa_{\mathrm F}}
\newcommand{\cP}{\mathcal{P}}
\newcommand{\cQ}{\mathcal{Q}}
\newcommand{\cL}{\mathcal{L}}
\newcommand{\rundir}{results/2026-11-24}
\newcolumntype{Y}{>{\raggedright\arraybackslash}X}
\newcolumntype{P}[1]{>{\raggedright\arraybackslash}p{#1}}

\theoremstyle{definition}
\newtheorem{example}{Example}[section]
\newtheorem{exercise}{Exercise}
\theoremstyle{remark}
\newtheorem{remark}{Remark}[section]
\numberwithin{equation}{section}

\title{Closure methods: a comprehensive tutorial\\[4pt]
\large Hierarchies, classical closures, plasma closures, learned closures,\\
and how to test them}
\author{Working note}
\date{24 November 2026}

\begin{document}
\maketitle

\noindent\emph{What this is.} An expanded version of the two earlier notes
\texttt{closure\_models\_tutorial.tex} (a high-level overview) and
\texttt{hammett\_perkins\_closure\_derivation.tex} (one derivation). Both are
kept unchanged in the parent directory, because the heat-flux manuscript
records their hashes as source notes. Everything in them is included here,
with derivations, worked examples, the closure families they left out
(turbulence, memory, magnetized and strongly coupled plasmas, radiation,
waves, learned closures), and a map onto the benchmark.

\medskip
\noindent\emph{Status.}
\begin{itemize}
  \item \textbf{Derived here and checked numerically.} The worked examples
    and every number in them: the fast--slow model
    (Section~\ref{sec:problem}), the cubic Langevin hierarchy
    (Section~\ref{sec:hier}), the BGK transport coefficients
    (Section~\ref{sec:asym}), the Hermite truncation
    (Section~\ref{sec:moment}), the Eddington factors
    (Section~\ref{sec:rad}), the Hammett--Perkins coefficients, the
    closure at general order and the error figures
    (Section~\ref{sec:collisionless}), and the comparison between the
    limits (Section~\ref{sec:between}). The script is \texttt{checks.py};
    its output is \texttt{\rundir/checks.json}, and
    Appendix~\ref{app:checks} lists what was compared with what.
  \item \textbf{Quoted from this repository.} Results of the closure notes
    and of the heat-flux manuscript are quoted with the conditions stated
    there (Section~\ref{sec:repo}). They are not re-proved here, and no
    proof in those notes has been checked by a second reader.
  \item \textbf{Written from general knowledge.} The standard theory, the
    historical attributions, and the numerical constants of named models
    (turbulence-model constants, Braginskii coefficients, nonlocal
    transport fits, freezing thresholds). These were not compared line by
    line with the cited sources. The status of the bibliography is given at
    the head of the reference list.
\end{itemize}

\noindent\emph{How to read it.} Sections~\ref{sec:problem}
and~\ref{sec:hier} set up the problem. Sections~\ref{sec:asym}
to~\ref{sec:memory} cover the four classical ways to close. Sections
\ref{sec:plasma} to~\ref{sec:waves} cover plasmas and the two other TEMPEST
settings. Sections~\ref{sec:learned} to~\ref{sec:repo} cover learned
closures, testing, and this project. A reader who wants only the argument
can read Sections~\ref{sec:problem}, \ref{sec:asym}, \ref{sec:memory},
\ref{sec:collisionless}, \ref{sec:between} and~\ref{sec:judge}.
Section~\ref{sec:summary} has summary tables, exercises with answers, and a
list of symbols.

\tableofcontents

\section{What a closure problem is}
\label{sec:problem}

\subsection{The abstract statement}

Suppose a system has a full state $X$ obeying $\dot X = F(X)$, and we only
want to track a reduced state $x = \cP X$: an average, a few moments, a
coarse-grained field. Applying $\cP$ to the dynamics gives
\begin{equation}
  \dot x = \cP F(X).
  \label{eq:unclosed}
\end{equation}
This equation is exact. Its right-hand side depends on the full state $X$,
not only on $x$, and that is the closure problem. A \emph{closure} is a rule
that replaces $\cP F(X)$ by something computable from $x$ alone, possibly
from its gradients or its history:
\begin{equation}
  \dot x = f(x) + \underbrace{\bigl[\cP F(X) - f(x)\bigr]}_{\text{closure defect } d} .
  \label{eq:defect}
\end{equation}
The bracket is what the closure gets wrong. It is a function of the full
state, so it can be measured only where the full state is available: on
reference trajectories or reference ensembles.

Two different objects are called ``the reduced equation'', and it helps to
keep them apart. Equation~\eqref{eq:unclosed} is an identity obtained by
averaging, filtering, taking moments or projecting. The closed model
$\dot x = f(x)$ is an approximation. The modelling enters only in the step
from the first to the second.

\subsection{Two ways to measure the error}

Let $X(t)$ be a reference solution, and let $x_c(t)$ solve the closed model
from the same reduced initial data.
\begin{itemize}
  \item The \textbf{a priori} error is the defect on the reference solution,
    $d(t) = \cP F(X(t)) - f(\cP X(t))$.
  \item The \textbf{a posteriori} (embedded) error is
    $e(t) = x_c(t) - \cP X(t)$.
\end{itemize}
Subtracting the two evolution equations gives
$\dot e = f(x_c) - f(\cP X) - d(t)$. If $f$ is Lipschitz with constant $L$
then by Gr\"onwall's inequality
\begin{equation}
  |e(t)| \le \int_0^t e^{L(t-s)}\,|d(s)|\,\dd s .
  \label{eq:gronwall}
\end{equation}
This inequality says three things. A small defect controls the embedded
error for times of order $1/L$ and no longer, unless the closed model is
contractive. The defect that matters is the one along the reference
trajectory, while a closed model run for a long time visits states the
reference never produced. And the bound is one-directional: a closure can
have a large pointwise defect and still give good embedded statistics, as
eddy-viscosity models do (Section~\ref{sec:turb}). These are the reasons the
benchmark requires an embedded-solver test next to every a priori test.

\subsection{Five items that describe any scale transition}

Every transition in this project is described by five items. Later sections
state them for each closure.
\begin{center}
\renewcommand{\arraystretch}{1.25}
\begin{tabularx}{\textwidth}{@{}P{3.4cm}Y@{}}
\toprule
Exact hierarchy & The chain of exact equations in which each retained
  quantity is coupled to one that is not retained. \\
Limiting process & The parameter that is sent to zero or infinity, and the
  scaling of time and space that goes with it. \\
Retained state & What the reduced model evolves. \\
Discarded information & What $\cP$ removes: velocity-space structure,
  correlations, fluctuations, phases, fast fields, history. \\
Closure defect & The bracket in \eqref{eq:defect}, with a statement of how
  it is measured. \\
\bottomrule
\end{tabularx}
\end{center}

\subsection{Where the problem comes from}

The problem has two causes, and usually both are present.
\begin{itemize}
  \item \textbf{Nonlinearity.} The average of a product is not the product of
    averages. Averaging a quadratic term produces a second moment, whose
    equation contains a third.
  \item \textbf{Transport.} Streaming or advection couples each retained
    quantity to one that is not retained. Multiplying
    $\bv\cdot\nabla f$ by $\bv^k$ and integrating always produces a moment
    of order $k+1$, even when the equation is linear.
\end{itemize}
A third situation is the elimination of whole degrees of freedom: fast
variables, or fields carried by another species. It is the same problem with
a different projection, and Section~\ref{sec:memory} treats it in general.

\subsection{A worked example with every feature}
\label{sec:fastslow}

\begin{example}[A slow variable driven by a fast one]
\label{ex:fastslow}
Let
\begin{equation}
  \dot x = a\,x + b\,y, \qquad
  \varepsilon\,\dot y = c\,x - y ,
  \label{eq:fastslow}
\end{equation}
with $\varepsilon>0$ small. Retain $x$ and discard $y$. The second equation
can be solved exactly,
\[
  y(t) = e^{-t/\varepsilon}\,y(0)
       + \frac{c}{\varepsilon}\int_0^t e^{-s/\varepsilon}\,x(t-s)\,\dd s ,
\]
so the exact reduced equation is
\begin{equation}
  \dot x(t) = a\,x(t)
   + \frac{bc}{\varepsilon}\int_0^t e^{-s/\varepsilon}\,x(t-s)\,\dd s
   + b\,e^{-t/\varepsilon}\,y(0) .
  \label{eq:fastslow-exact}
\end{equation}
The three terms are an instantaneous term, a \emph{memory} term, and a term
that depends on the discarded initial data. Nothing has been approximated.

Three closures follow from expanding the memory integral for slowly varying
$x$, using $x(t-s) = x(t) - s\,\dot x(t) + \dots$:
\begin{align}
  \text{zeroth order (Markov):}\quad & \dot x = (a+bc)\,x, \nonumber\\
  \text{first order:}\quad & \dot x = (a+bc)(1-\varepsilon bc)\,x, \\
  \text{first order, rational:}\quad
     & (1+\varepsilon bc)\,\dot x = (a+bc)\,x . \nonumber
\end{align}
The second line substitutes the zeroth-order equation for $\dot x$ in the
correction. The third keeps $\dot x$ as an unknown. They agree to
$O(\varepsilon)$ and differ at $O(\varepsilon^2)$.

For $a=0$, $b=1$, $c=-1$ the exact slow eigenvalue is
$\lambda = (-1+\sqrt{1-4\varepsilon})/(2\varepsilon)
= -1-\varepsilon-2\varepsilon^2-\dots$:
\begin{center}
\begin{tabular}{rrrrr}
\toprule
$\varepsilon$ & exact & Markov & first order & rational \\
\midrule
0.2   & $-1.3820$ & $-1$ & $-1.2000$ & $-1.2500$ \\
0.1   & $-1.1270$ & $-1$ & $-1.1000$ & $-1.1111$ \\
0.05  & $-1.0557$ & $-1$ & $-1.0500$ & $-1.0526$ \\
0.025 & $-1.0263$ & $-1$ & $-1.0250$ & $-1.0256$ \\
\bottomrule
\end{tabular}
\end{center}
The Markov error is $O(\varepsilon)$ and the first-order errors are
$O(\varepsilon^2)$.
\end{example}

The example contains the whole subject in miniature.
\begin{itemize}
  \item The exact reduced dynamics has memory. A local closure is obtained
    only by assuming that the memory kernel decays faster than $x$ changes.
    That is a statement about a small parameter, here $\varepsilon$, in
    kinetic theory the Knudsen number.
  \item The closure is \emph{asymptotic}: the kernel
    $\varepsilon^{-1}e^{-s/\varepsilon}$ tends to a delta function, and the
    corrections are a series in $\varepsilon$.
  \item The discarded initial data matter for a time of order $\varepsilon$
    (an initial layer). If they are unknown they act as noise.
  \item The same order of accuracy can be reached by formulas that behave
    very differently when $\varepsilon$ is not small. The polynomial form
    changes sign at $\varepsilon bc = 1$. The rational form does not. The
    same distinction separates the Burnett equations from their
    regularizations (Section~\ref{sec:burnett}).
  \item Keeping $y$ as a variable removes the closure problem at the cost of
    a larger, stiffer system. Extended-moment models make this trade
    (Section~\ref{sec:moment}).
\end{itemize}

\subsection{The kinds of closure}
\label{sec:kinds}

Closures differ in what the unclosed quantity is allowed to depend on.
\begin{center}
\renewcommand{\arraystretch}{1.25}
\begin{tabularx}{\textwidth}{@{}P{3.3cm}P{5.3cm}Y@{}}
\toprule
Kind & Form & Examples \\
\midrule
Algebraic & $u = g(x)$ at the same point and time
  & equation of state; Euler ($\boldsymbol\pi=\bq=0$); double-adiabatic
    laws; Eddington factor \\
Gradient & $u = -D(x)\,\nabla x$
  & Navier--Stokes--Fourier; eddy viscosity; Braginskii; diffusion
    approximation \\
Nonlocal in space & $u(\bx) = \int W(\bx,\bx')\,x(\bx')\,\dd\bx'$
  & Hammett--Perkins; nonlocal electron heat flux \\
Memory & $u(t) = \int_0^t K(s)\,x(t-s)\,\dd s$
  & Mori--Zwanzig; viscoelastic models \\
Extended state & evolve $u$ with its own closed equation
  & Grad; Reynolds-stress transport; $k$--$\varepsilon$; Cattaneo \\
Stochastic & $u = g(x) + \text{noise}$
  & Langevin models; stochastic subgrid forcing \\
\bottomrule
\end{tabularx}
\end{center}
These are not exclusive. An extended-state model is a memory closure with an
exponential kernel (Section~\ref{sec:embedding}), and for linear problems
every row is a choice of how to approximate one function of wavenumber and
frequency (Section~\ref{sec:between}).

\section{Where the hierarchies come from}
\label{sec:hier}

Five hierarchies recur in this project. Each is exact, and each has the
same shape: the equation at level $s$ contains the unknown of level $s+1$.

\subsection{Particles to kinetic: BBGKY}

Take $N$ identical particles of mass $m$ with pair potential $\phi$, and let
$F_N$ be a symmetric probability density on phase space obeying the
Liouville equation. The $s$-particle marginal is
$f_s = \int F_N\,\dd z_{s+1}\cdots\dd z_N$ with $z_i = (\bx_i,\bv_i)$.
Integrating the Liouville equation gives the BBGKY hierarchy. Its first
equation is
\begin{equation}
  \partial_t f_1 + \bv\cdot\nabla_{\bx} f_1
  = \frac{N-1}{m}\int \nabla\phi(\bx-\bx')\cdot\nabla_{\bv}
    f_2(\bx,\bv,\bx',\bv',t)\,\dd\bx'\,\dd\bv' ,
  \label{eq:bbgky1}
\end{equation}
and in general
\begin{equation}
  \partial_t f_s = \cL_s f_s + \mathcal{C}_{s,s+1} f_{s+1},
  \qquad s = 1,\dots,N-1,
\end{equation}
where $\cL_s$ is the dynamics of $s$ particles among themselves and
$\mathcal{C}_{s,s+1}$ is their interaction with one more. Here the cause is
the interaction, which is a two-body quantity.

The useful variables are the connected correlations (cumulants),
\[
  f_2 = f_1 f_1 + g_2, \qquad
  f_3 = f_1 f_1 f_1 + \textstyle\sum f_1\,g_2 + g_3, \qquad\dots
\]
A kinetic equation is a closure of \eqref{eq:bbgky1}. Which one depends on
the scaling.
\begin{center}
\renewcommand{\arraystretch}{1.25}
\begin{tabularx}{\textwidth}{@{}P{3.1cm}P{4.2cm}P{3.0cm}Y@{}}
\toprule
Scaling & Small parameter & Closure & Kinetic equation \\
\midrule
Low density (Boltzmann--Grad)
  & hard spheres of diameter $\epsilon$, $N\epsilon^{d-1}$ fixed
  & $g_2\to0$ for particles about to collide
  & Boltzmann \\
Mean field
  & $\phi\to\phi/N$, long range
  & $g_2 = O(1/N)$ neglected
  & Vlasov \\
Weak coupling
  & weak potential, many grazing encounters
  & $g_3$ neglected, $g_2$ slaved to $f_1$
  & Landau \\
Plasma
  & $1/\Lambda$, particles per Debye volume
  & as above, with collective screening
  & Vlasov at leading order, Lenard--Balescu at next \\
\bottomrule
\end{tabularx}
\end{center}
In the first row the assumption is one-sided: particles are uncorrelated
\emph{before} they collide, and correlated afterwards. That asymmetry is
where irreversibility enters. Lanford proved the first row for short times
\cite{Lanford75}; Deng, Hani and Ma extended it to the full lifespan of the
Boltzmann solution \cite{DHM24}, by propagating a cumulant ansatz across
time layers. Their organization of $g_2, g_3,\dots$ by collision history is
the structure the mathematical closure program in this repository tries to
turn into measurable diagnostics.

\subsection{Kinetic to fluid: the moment hierarchy}
\label{sec:momenthier}

A dilute gas is described by $f(\bx,\bv,t)$ obeying
\begin{equation}
  \partial_t f + \bv\cdot\nabla f + \frac{\bF}{m}\cdot\nabla_{\bv} f = C[f],
  \label{eq:kinetic}
\end{equation}
with $C$ the collision operator and $\bF$ a force. Define density
$n=\int f\,\dd^3v$, flow $n\bu=\int\bv f\,\dd^3v$, and temperature
$\tfrac32 nT=\int\tfrac12 m|\bv-\bu|^2f\,\dd^3v$ (with $k_B=1$). Because
collisions conserve particles, momentum and energy, the first three moment
equations contain no collision term:
\begin{align}
  \partial_t n + \nabla\cdot(n\bu) &= 0, \label{eq:mass}\\
  mn\,\frac{D\bu}{Dt} &= -\nabla p - \nabla\cdot\boldsymbol{\pi} + n\bF,
     \label{eq:mom}\\
  \frac{3}{2}\,n\,\frac{DT}{Dt} + p\,\nabla\cdot\bu
     &= -\nabla\cdot\bq - \boldsymbol{\pi}:\nabla\bu , \label{eq:energy}
\end{align}
with $p=nT$, $D/Dt = \partial_t+\bu\cdot\nabla$, and
\[
  \boldsymbol\pi = m\!\int\!\bigl(\bw\bw-\tfrac13|\bw|^2\mathbf I\bigr)f\,\dd^3v,
  \qquad
  \bq = \int \tfrac12 m|\bw|^2\,\bw\,f\,\dd^3v,
  \qquad \bw=\bv-\bu .
\]
These are exact and not closed: $\boldsymbol\pi$ and $\bq$ are higher
moments. Their equations contain still higher moments, and also moments of
$C[f]$, which are not zero beyond the conserved ones. Maxwell wrote the
hierarchy down in 1867.

In one space and one velocity dimension the structure is plain. With
$M_k=\int v^kf\,\dd v$ and a force-free, collisionless equation,
\begin{equation}
  \partial_t M_k + \partial_x M_{k+1} = 0, \qquad k=0,1,2,\dots
  \label{eq:rawhier}
\end{equation}
Retaining $M_0,\dots,M_N$ leaves $M_{N+1}$ undetermined. Here the cause is
transport alone; the equation is linear.

\subsection{Fluid to mean flow: Reynolds averaging and filtering}

For an incompressible fluid split $u_i=\bar u_i+u_i'$, where the bar is an
ensemble average, and average the Navier--Stokes equations (Reynolds, 1895):
\begin{equation}
  \partial_t \bar u_i + \bar u_j\,\partial_j \bar u_i
  = -\frac{1}{\rho}\,\partial_i \bar p + \nu\,\nabla^2 \bar u_i
    - \partial_j \avg{u_i' u_j'} .
  \label{eq:rans}
\end{equation}
The Reynolds stress $R_{ij}=\avg{u_i'u_j'}$ is not a function of the mean
flow. Its equation contains triple correlations and pressure--velocity
correlations, and so on. The cause is the quadratic nonlinearity.

Large-eddy simulation replaces the average by a spatial filter of width
$\Delta$. The filtered equations have the same form with the
\emph{subgrid stress}
\begin{equation}
  \tau_{ij} = \overline{u_iu_j}-\bar u_i\,\bar u_j
  \label{eq:sgs}
\end{equation}
in place of $R_{ij}$. Two differences matter later. A filter is not a
projection in general ($\bar{\bar u}\ne\bar u$), and $\tau_{ij}$ fluctuates
in space and time, so an LES closure is asked for a field and not a
statistic.

\subsection{A hierarchy of statistical moments, solved}
\label{sec:langevin}

The same structure appears whenever a nonlinear system is described by
statistics.

\begin{example}[Cumulant-discard closure for a nonlinear Langevin equation]
\label{ex:langevin}
Let $\dd x=-x^3\,\dd t+\sigma\,\dd W$ and $m_k=\mathbb E\,x^k$. It\^o's
formula gives
\begin{equation}
  \dot m_k = -k\,m_{k+2} + \tfrac12\sigma^2\,k(k-1)\,m_{k-2} .
\end{equation}
The equation for $m_2$ contains $m_4$, that for $m_4$ contains $m_6$. The
stationary density is $p\propto\exp(-x^4/2\sigma^2)$, so the exact answers
are
\[
  m_2=\sqrt2\,\frac{\Gamma(3/4)}{\Gamma(1/4)}\,\sigma = 0.4780\,\sigma,
  \qquad m_4=\tfrac12\sigma^2 .
\]
\emph{Gaussian closure.} Set the fourth cumulant to zero, $m_4=3m_2^2$.
Then $\dot m_2=-6m_2^2+\sigma^2$, with stationary value
$m_2=\sigma/\sqrt6=0.4082\,\sigma$, which is $14.6\%$ low.

\emph{Next order.} Keep $m_2$ and $m_4$ and set the sixth cumulant to
zero, $m_6=15m_2m_4-30m_2^3$. The stationary equations give
$m_4=\sigma^2/2$ exactly and $m_2=\sigma/\sqrt5=0.4472\,\sigma$, which is
$6.4\%$ low.
\end{example}

Three remarks carry over to turbulence and to plasmas. The closure assumes
a \emph{shape} for the distribution (Gaussian, or Gaussian plus a fourth
cumulant), which is closure by ansatz (Section~\ref{sec:moment}). No small
parameter makes the discarded cumulant small, so the error is not
controlled and improvement with order is not guaranteed. And a set of
closed moments need not be the moments of any probability density; when it
is not, the closure is not \emph{realizable} (Section~\ref{sec:realizability}).
Discarding cumulants above the second is called the quasi-normal
approximation in turbulence, and in that setting it does produce negative
energy spectra (Section~\ref{sec:statclosure}).

\subsection{Waves to wave kinetics}

For weakly nonlinear dispersive waves, for example the cubic nonlinear
Schr\"odinger equation
$i\partial_tu+\Delta u=\lambda|u|^2u$ with random initial data, the retained
quantity is the spectrum $n_k=\mathbb E|\hat u_k|^2$. Its equation contains
fourth-order correlations of Fourier amplitudes, whose equation contains
sixth-order ones. The closure is the random-phase (near-Gaussian)
approximation, the analogue of molecular chaos, and the result is the wave
kinetic equation (Section~\ref{sec:waves}).

\subsection{The common pattern}

\begin{center}
\renewcommand{\arraystretch}{1.3}
\begin{tabularx}{\textwidth}{@{}P{3.0cm}P{2.7cm}P{2.9cm}Y@{}}
\toprule
Problem & Retained state & Unclosed quantity & Discarded information \\
\midrule
Particles to kinetic & $f_1$ & $f_2$ (or $g_2$)
  & correlations between particles, collision history \\
Kinetic to fluid & $n$, $\bu$, $T$ & $\boldsymbol\pi$, $\bq$
  & shape of $f$ in velocity \\
Turbulence (RANS, LES) & $\bar{\bu}$, $\bar p$
  & $\avg{u_i'u_j'}$ or $\tau_{ij}$
  & fluctuations below the average or filter \\
Waves to wave kinetic & $n_k$ & fourth-order correlations
  & phases, higher spectral cumulants \\
Radiation transport & energy density, flux & pressure tensor
  & angular distribution of the radiation \\
Dusty plasma particles & grain positions and velocities
  & force on each grain
  & ions, electrons, wakes and fields \\
3D to 1D (supernova) & spherical averages
  & turbulent stresses and fluxes
  & nonradial motion, shock corrugation, history \\
\bottomrule
\end{tabularx}
\end{center}

\section{Asymptotic closure}
\label{sec:asym}

If a small parameter makes the discarded variables relax quickly to a state
determined by the retained ones, the closure can be \emph{derived}, with an
error estimate. This is the only kind of closure that comes with one.

\subsection{The small parameter and two expansions}

For a gas the parameter is the Knudsen number
\begin{equation}
  \Kn = \frac{\lambda}{L},
\end{equation}
the mean free path over the gradient length. In units of the macroscopic
length and time the Boltzmann equation reads
$\partial_tf+\bv\cdot\nabla f=\Kn^{-1}C[f]$. As $\Kn\to0$ the collision
term forces $C[f]=0$, whose solutions are the local Maxwellians
\[
  f_M(n,\bu,T) = n\Bigl(\frac{m}{2\pi T}\Bigr)^{3/2}
  \exp\Bigl(-\frac{m|\bv-\bu|^2}{2T}\Bigr).
\]
Two expansions organize the corrections.
\begin{itemize}
  \item \textbf{Hilbert (1912)} expands both $f$ and the fluid fields in
    powers of $\Kn$. Every order solves linearized Euler equations forced by
    the previous orders. The result is a series solution, not a new set of
    equations.
  \item \textbf{Chapman--Enskog (1916--17)} expands $f$ but not the fluid
    fields:
    $f=f_M(n,\bu,T)\,(1+\Kn\,\phi_1+\Kn^2\phi_2+\dots)$, with each
    $\phi_j$ a function of the fields and their gradients and carrying no
    density, momentum or energy. The time derivative is expanded instead.
    The result is a sequence of closed equations for $n$, $\bu$, $T$: Euler,
    Navier--Stokes--Fourier, Burnett.
\end{itemize}

\subsection{First order: Navier--Stokes--Fourier}
\label{sec:ce}

Let $\Lambda$ be the linearized collision operator about $f_M$, written so
that $C[f_M(1+\phi)] = -f_M\,\Lambda\phi+O(\phi^2)$. It is self-adjoint and
non-negative in the Maxwellian-weighted inner product
$\avg{\phi,\psi}=n^{-1}\int\phi\psi f_M\,\dd^3v$, and its kernel is spanned
by the collision invariants $1$, $\bv$, $|\bv|^2$.

At order $\Kn^0$ the kinetic equation gives
\begin{equation}
  \Lambda\phi_1 = -\bigl(\partial_t+\bv\cdot\nabla\bigr)\ln f_M ,
  \label{eq:ce1}
\end{equation}
where the time derivatives on the right are replaced by the Euler
equations. Writing $\bw=\bv-\bu$, a direct computation gives
\begin{equation}
  \bigl(\partial_t+\bv\cdot\nabla\bigr)\ln f_M
  = \Bigl(\frac{m|\bw|^2}{2T}-\frac52\Bigr)\bw\cdot\nabla\ln T
  + \frac{m}{T}\Bigl(\bw\bw-\tfrac13|\bw|^2\mathbf I\Bigr):\nabla\bu .
  \label{eq:ce-source}
\end{equation}
Three points deserve attention.
\begin{enumerate}
  \item \emph{Solvability.} Equation \eqref{eq:ce1} can be solved only if
    its right-hand side is orthogonal to the kernel of $\Lambda$. The Euler
    equations are exactly that condition. This is the Fredholm alternative,
    and it is where the fluid equations come from at every order.
  \item \emph{The source has two pieces,} a vector multiplying $\nabla T$
    and a traceless tensor multiplying $\nabla\bu$. By isotropy of
    $\Lambda$ the first produces only a heat flux and the second only a
    stress. No density-gradient term survives.
  \item \emph{The transport coefficients are quadratic forms.} With
    dimensionless $\bc=\bw/\vt$, $\vt=\sqrt{T/m}$, and
    \[
      \mathbf A=\bigl(\tfrac12|\bc|^2-\tfrac52\bigr)\bc,\qquad
      \mathbf B=\bc\bc-\tfrac13|\bc|^2\mathbf I,
    \]
    the fluxes are
    \begin{equation}
      \boldsymbol{\pi} = -\mu\left(\nabla\bu + \nabla\bu^{\mathsf T}
         - \tfrac{2}{3}(\nabla\cdot\bu)\,\mathbf{I}\right),
      \qquad
      \bq = -\kappa\,\nabla T ,
      \label{eq:NSF}
    \end{equation}
    \begin{equation}
      \kappa = \frac{nT}{m}\cdot\frac13\avg{\mathbf A,\Lambda^{-1}\mathbf A},
      \qquad
      \mu = nT\cdot\frac1{10}\avg{\mathbf B,\Lambda^{-1}\mathbf B} .
      \label{eq:transport-forms}
    \end{equation}
\end{enumerate}
Because $\Lambda$ is non-negative, each coefficient is the maximum of a
concave quadratic functional,
\begin{equation}
  \avg{\mathbf A,\Lambda^{-1}\mathbf A}
  = \max_{\boldsymbol\varphi}\bigl\{2\avg{\mathbf A,\boldsymbol\varphi}
      -\avg{\boldsymbol\varphi,\Lambda\boldsymbol\varphi}\bigr\},
  \label{eq:variational}
\end{equation}
so any trial function gives a lower bound. This is how transport
coefficients have been computed since Chapman and Enskog (polynomial trial
functions, the Sonine expansion), and it is the starting point of the
two-sided bounds proved in this repository (Section~\ref{sec:repo}).

\begin{example}[BGK]
\label{ex:bgk}
The BGK model replaces collisions by relaxation to the local Maxwellian at
rate $\nu$, $C[f]=\nu(f_M-f)$ \cite{BGK54}. Then $\Lambda=\nu$ on the
complement of the collision invariants, and \eqref{eq:transport-forms}
needs only Gaussian moments:
$\tfrac13\avg{\mathbf A,\mathbf A}=\tfrac52$ and
$\tfrac1{10}\avg{\mathbf B,\mathbf B}=1$. Hence
\begin{equation}
  \kappa=\frac52\,\frac{nT}{m\nu},\qquad \mu=\frac{nT}{\nu},\qquad
  \Pr=\frac{c_p\,\mu}{\kappa}=1 \quad(c_p=\tfrac52/m).
\end{equation}
A monatomic gas has $\Pr\approx\tfrac23$. With one relaxation rate BGK
cannot match the viscosity and the conductivity together, which is a first
instance of a theme of Section~\ref{sec:between}: a model collision
operator fitted to one transport coefficient does not reproduce the others.

In one velocity dimension, with the heat flux defined as the third central
moment (the convention of Section~\ref{sec:collisionless}), the same
calculation gives $\kF=3\,nT/(m\nu)$.
\end{example}

For hard spheres of diameter $d$ the first Sonine approximation gives
$\mu=\tfrac{5}{16d^2}\sqrt{mT/\pi}$ and $\Pr=\tfrac23$; higher
approximations change $\mu$ and $\kappa$ by about two percent
\cite{ChapmanCowling}. In both cases
$\mu\sim mn\,\vt\lambda$ and $\kappa\sim n\,\vt\lambda$: a mean free path
times a thermal speed.

\subsection{Higher order, and why it fails}
\label{sec:burnett}

The second order of the Chapman--Enskog expansion gives the Burnett
equations, with terms of second order in gradients. They are not an
improvement in practice. Bobylev showed that they are linearly unstable at
short wavelengths \cite{Bobylev82}.

The mechanism can be seen in one line. For a linear temperature
perturbation $\propto e^{ikx}$ write the heat flux as
$\hat q=-ik\,\chi(k)\,\hat T$, so that the temperature equation is
$\partial_t\hat T\propto-k^2\chi(k)\hat T$. Fourier's law is
$\chi=\kF$. The expansion continued in powers of the gradient gives
\begin{equation}
  \chi(k)=\kF\bigl(1-K\,k^2+O(k^4)\bigr),\qquad K>0,
  \label{eq:deficit}
\end{equation}
where the sign reflects that the true flux falls below Fourier's as
gradients steepen. Truncating after the $k^2$ term makes $\chi$ negative for
$k^2>1/K$, where the temperature mode grows at a rate proportional to
$k^4$: an ill-posed equation. The rational form
\begin{equation}
  \chi(k)=\frac{\kF}{1+K\,k^2}
  \label{eq:rational}
\end{equation}
has the same expansion to this order and is positive and bounded at every
$k$. This is the same distinction as in Example~\ref{ex:fastslow} between
the polynomial and the rational first-order closures. For conserving
collision operators in one velocity dimension, \eqref{eq:rational} is a
proved lower bound on the exact response, and $\kF$ is an upper bound
(Section~\ref{sec:repo}).

Repairs of the Burnett equations take three routes: regularized moment
systems such as R13 \cite{StruchtrupTorrilhon03,Struchtrup05}; exact
summation of the expansion in linear problems, and invariant-manifold
constructions \cite{GorbanKarlin14,KogelbauerKarlin24}; and abandoning the
expansion for a moment closure (Section~\ref{sec:moment}).

\subsection{What is proved}

The formal expansions are backed by theorems in several settings. Caflisch
justified the compressible Euler limit for smooth solutions
\cite{Caflisch80}. Bardos, Golse and Levermore set out the program for
incompressible limits \cite{BGL91}, and Golse and Saint-Raymond carried it
through for the incompressible Navier--Stokes limit in the setting of
renormalized solutions \cite{GSR04}. Deng, Hani and Ma connect hard-sphere
dynamics through the Boltzmann equation to the fluid equations
\cite{DHM25}. These are the theorem-faithful settings that Track~A of the
benchmark uses.

\subsection{Why neutral gases are the easy case}
\label{sec:easy}

Collisions between neutral molecules are short-range and frequent. Three
facts follow, and each fails for plasmas (Section~\ref{sec:plasma}).
\begin{enumerate}
  \item \textbf{Fast local relaxation.} Collisions drive $f$ to a local
    Maxwellian in a few collision times, so the kinetic state is determined
    by $n$, $\bu$, $T$ and their gradients.
  \item \textbf{One small parameter.} All molecules have comparable mean
    free paths, so a single $\Kn$ controls the expansion.
  \item \textbf{Locality.} Fluxes depend on gradients at the same point and
    time.
\end{enumerate}

\section{Closure by an ansatz for the distribution}
\label{sec:moment}

When no small parameter is available, or when a model must stay valid away
from equilibrium, the alternative is to retain more moments and assume a
shape for $f$ with as many free parameters as retained moments. The missing
moment is then computed from that shape. These closures give evolution
equations for the fluxes, so the fluxes have finite relaxation times and
signals have finite speeds.

\subsection{The linear model problem in a Hermite basis}
\label{sec:hermite}

One linear problem is used repeatedly below. Take one space and one
velocity dimension, units with $n_0=m=\vt=1$, the Maxwellian
$M(v)=(2\pi)^{-1/2}e^{-v^2/2}$, and write $f=M(1+h)$. The linearized
kinetic equation with acceleration $a(x,t)$ is
\begin{equation}
  \partial_th+v\,\partial_xh+\Lambda h=a\,v .
  \label{eq:linkin}
\end{equation}
Expand $h=\sum_{n\ge0}a_n(x,t)\,\psi_n(v)$ in the orthonormal Hermite
functions $\psi_n=\He_n/\sqrt{n!}$, where $\He_0=1$, $\He_1=v$,
$\He_2=v^2-1$, $\He_3=v^3-3v$. The recurrence
$v\,\psi_n=\sqrt{n+1}\,\psi_{n+1}+\sqrt n\,\psi_{n-1}$ turns
\eqref{eq:linkin} into
\begin{equation}
  \partial_ta_n+\partial_x\bigl(\sqrt{n+1}\,a_{n+1}+\sqrt n\,a_{n-1}\bigr)
  +(\Lambda a)_n=a\,\delta_{n1} .
  \label{eq:hermitehier}
\end{equation}
The coefficients are the fluid moments: $a_0$ is the density perturbation,
$a_1$ the velocity, $\sqrt2\,a_2$ the temperature, $\sqrt6\,a_3$ the heat
flux (third central moment), $\sqrt{24}\,a_4$ the deviation of the fourth
moment from its Maxwellian value. For BGK, $\Lambda$ is diagonal with
entries $0,0,0,\nu,\nu,\dots$; for the Dougherty (Lenard--Bernstein)
operator \cite{LenardBernstein58,Dougherty64} the entries are
$0,0,0,3\nu,4\nu,\dots$.

Streaming is a tridiagonal (Jacobi) matrix in this basis. It couples each
coefficient to its two neighbours and nothing else, and all closure
questions for the linear problem are questions about what to do at the last
retained row.

\subsection{Truncation}

The simplest choice is $a_{N+1}=0$. For a linear problem this is Grad's
closure. Its content is exact and worth knowing.

\begin{example}[Truncation is a discrete-velocity model]
With $a_{N+1}=0$ the collisionless system is
$\partial_t\mathbf a+J_N\,\partial_x\mathbf a=0$, where $J_N$ is the
$(N+1)\times(N+1)$ Jacobi matrix with off-diagonal entries
$\sqrt1,\dots,\sqrt N$. Its eigenvalues are the roots of $\He_{N+1}$, the
Gauss--Hermite quadrature nodes. (Checked for $N=8$: the eigenvalues and
the nodes agree to $3\times10^{-15}$, and the largest is $4.51$.) The
truncated moment system is therefore equivalent to a kinetic equation with
$N+1$ discrete velocities and Gauss--Hermite weights.
\end{example}

Two consequences follow. The truncated system is hyperbolic with real wave
speeds, so it is well posed. And without collisions it cannot damp: a
density perturbation at wavenumber $k$ is a finite sum of oscillations
$e^{-ikv_jt}$, which decays for a while and then returns. The recurrence
time is of order $\sqrt N/(|k|\vt)$. Landau damping
(Section~\ref{sec:collisionless}) is lost, however many moments are kept,
unless something absorbs the flow of perturbation energy to high $n$. That
absorber is either collisions, which act strongly at high $n$, or a
closure that models the outflow.

\subsection{Grad's method}

Grad \cite{Grad49} expands the nonlinear distribution as a local Maxwellian
times a polynomial,
\[
  f=f_M(n,\bu,T)\Bigl[1+\frac{\boldsymbol\pi:\bc\bc}{2p}
     +\frac{2}{5}\,\frac{\bq\cdot\bc}{p\,\vt}
       \Bigl(\frac{|\bc|^2}{2}-\frac52\Bigr)\Bigr],
\]
which has thirteen parameters $(n,\bu,T,\boldsymbol\pi,\bq)$, and takes
moments of the Boltzmann equation to get thirteen evolution equations. The
unclosed higher moments are evaluated with this $f$. The standard counts
are 5 (Euler), 10 (Gaussian), 13, 14, 20, 26 and 35 moments.

Grad's equations contain Navier--Stokes--Fourier as the limit in which
$\boldsymbol\pi$ and $\bq$ relax quickly. They have two known defects. The
ansatz is not positive, and the equations lose hyperbolicity at moderate
departures from equilibrium, which appears as unphysical sub-shocks.
Regularized (R13) and globally hyperbolic variants address these
\cite{StruchtrupTorrilhon03,Struchtrup05}.

\subsection{Maximum entropy}

Levermore \cite{Levermore96} chooses the distribution that maximizes the
entropy $-\int(f\ln f-f)\,\dd v$ subject to the retained moments
$\int m_i(\bv)f\,\dd v=\rho_i$. The maximizer is
\begin{equation}
  f=\exp\Bigl(\sum_i\alpha_i\,m_i(\bv)\Bigr),
\end{equation}
with multipliers $\alpha_i$ determined by the moments. The resulting system
has three structural properties by construction: it is symmetric hyperbolic
in the variables $\alpha$, it dissipates an entropy, and the ansatz is
positive, so its moments are realizable.

The costs are also structural. The map from moments to multipliers has a
closed form only for the 5-moment (Euler) and 10-moment (Gaussian) systems,
so each flux evaluation is an optimization problem. And beyond ten moments
the set of moments for which the problem is solvable is not what one would
hope: Junk showed that for the simplest higher system equilibrium lies on
the boundary of that set, and the fluxes are singular there \cite{Junk98}.
A closed-form example in another setting, the M1 closure for radiation, is
worked in Section~\ref{sec:rad}.

\subsection{Realizability and hyperbolicity}
\label{sec:realizability}

A vector of moments is \emph{realizable} if it comes from some non-negative
distribution. In one velocity dimension this is the Hamburger moment
problem: $(M_0,\dots,M_{2r})$ is realizable when the Hankel matrices
$(M_{i+j})_{i,j=0}^{r}$ are positive semidefinite. The first conditions are
$M_0\ge0$, $M_0M_2-M_1^2\ge0$ (non-negative temperature), and a lower bound
on the fourth moment in terms of the lower ones. A closure can leave the
realizable set in two ways: its formula for the unclosed moment can be
inconsistent with the retained ones, or the closed dynamics can carry
realizable data out of the set.

\emph{Hyperbolicity} is the requirement that the flux Jacobian of the
closed first-order system have real eigenvalues and a full set of
eigenvectors. Without it the initial-value problem is ill posed, and no
amount of accuracy in the closure term compensates. For learned closures
this is the property most often lost when a network is fitted to the
unclosed moment directly (Section~\ref{sec:learned}).

\subsection{Other ansatz closures}

\begin{itemize}
  \item \textbf{Gaussian (10-moment).} Retain the full pressure tensor and
    set the heat flux to zero. Used in plasma modelling where pressure
    anisotropy matters; the heat-flux closure added on top is often a
    relaxation term built on the Hammett--Perkins idea \cite{Wang15}.
  \item \textbf{Spectral truncations.} $P_N$ in radiation transport is
    truncation of a spherical-harmonic expansion; it is linear and can
    produce negative densities.
  \item \textbf{Quadrature-based methods.} Represent $f$ as a sum of delta
    functions (or kernels) whose weights and nodes are recovered from the
    retained moments \cite{McGraw97}. Realizable by construction; the
    inversion is ill conditioned at high order.
\end{itemize}

\begin{center}
\renewcommand{\arraystretch}{1.25}
\begin{tabularx}{\textwidth}{@{}P{2.9cm}P{3.5cm}YY@{}}
\toprule
Closure & Ansatz & Guarantees & Known failures \\
\midrule
Truncation / Grad & Maxwellian times polynomial
  & linear problem hyperbolic; recovers NSF
  & not positive; nonlinear loss of hyperbolicity; no collisionless damping \\
Maximum entropy & exponential of polynomial
  & hyperbolic, entropy, realizable
  & costly; singular fluxes beyond 10 moments \\
Quadrature & sum of deltas
  & realizable
  & ill-conditioned inversion; weakly hyperbolic variants \\
Gaussian & anisotropic Maxwellian
  & hyperbolic, realizable, closed form
  & no heat flux \\
\bottomrule
\end{tabularx}
\end{center}

\section{Turbulence closures}
\label{sec:turb}

Turbulence is the case with no small parameter: eddies of all sizes between
the mean-flow scale and the dissipation scale are present, and none is
separated from the scale being resolved. The closures here are
\emph{models}. They are included because the supernova profile needs them,
because most of the vocabulary for testing closures (a priori and a
posteriori tests, realizability) was developed for them, and because the
contrast with Section~\ref{sec:asym} is instructive. The standard reference
is Pope \cite{Pope}.

\subsection{Eddy viscosity}

Postulate the form of \eqref{eq:NSF} and model the coefficient. For the
Reynolds stress (Boussinesq, 1877),
\begin{equation}
  \avg{u_i'u_j'}-\tfrac23k\,\delta_{ij}=-2\nu_t\,\bar S_{ij},
  \qquad
  \bar S_{ij}=\tfrac12(\partial_i\bar u_j+\partial_j\bar u_i),
  \qquad k=\tfrac12\avg{u_i'u_i'} .
  \label{eq:boussinesq}
\end{equation}
The assumption is tensorial: the anisotropic stress is aligned with the
local mean strain. The eddy viscosity $\nu_t$ is then supplied at one of
several levels.
\begin{itemize}
  \item \textbf{Mixing length} (Prandtl, 1925):
    $\nu_t=\ell_m^2\,|\partial\bar u/\partial y|$ with a prescribed length
    $\ell_m$. This is the kinetic estimate $\nu\sim\vt\lambda$ with eddies
    in place of molecules. The analogy is not a derivation, because the
    ``mean free path'' of an eddy is comparable to the scale of the mean
    flow.
  \item \textbf{Two-equation models.} Transport two turbulence quantities
    and form $\nu_t$ from them. In the $k$--$\varepsilon$ model
    \cite{LaunderSpalding74},
    \begin{align}
      \nu_t &= C_\mu\,\frac{k^2}{\varepsilon}, \nonumber\\
      \frac{\bar Dk}{\bar Dt} &= P-\varepsilon
         +\nabla\cdot\Bigl[\Bigl(\nu+\frac{\nu_t}{\sigma_k}\Bigr)\nabla k\Bigr],
         \qquad P=2\nu_t\,\bar S_{ij}\bar S_{ij}, \\
      \frac{\bar D\varepsilon}{\bar Dt} &= C_{\varepsilon1}\frac{\varepsilon}{k}P
         -C_{\varepsilon2}\frac{\varepsilon^2}{k}
         +\nabla\cdot\Bigl[\Bigl(\nu+\frac{\nu_t}{\sigma_\varepsilon}\Bigr)
           \nabla\varepsilon\Bigr], \nonumber
    \end{align}
    with standard constants $C_\mu=0.09$, $C_{\varepsilon1}=1.44$,
    $C_{\varepsilon2}=1.92$, $\sigma_k=1.0$, $\sigma_\varepsilon=1.3$.
\end{itemize}
The constants are calibrated on canonical flows: $C_{\varepsilon2}$ from
the decay of grid turbulence, $C_\mu$ from the ratio of shear stress to $k$
in the logarithmic layer, and the others from requiring the logarithmic
law. The extra equations do not remove the closure problem. The $k$
equation is derived and then modelled term by term, and the $\varepsilon$
equation is largely constructed by analogy.

\emph{Realizability.} Equation \eqref{eq:boussinesq} gives
$\avg{u_1'^2}=\tfrac23k-2\nu_t\bar S_{11}$, which is negative when
$\bar S_{11}>k/(3\nu_t)$. A true Reynolds stress is a covariance and is
positive semidefinite. Realizable variants limit $C_\mu$ in strong strain.

\subsection{Second-moment closures}

Reynolds-stress transport models evolve all six components of $R_{ij}$,
\[
  \frac{\bar DR_{ij}}{\bar Dt}=P_{ij}+\Phi_{ij}-\varepsilon_{ij}+D_{ij} .
\]
Production $P_{ij}$ is closed. The pressure--strain redistribution
$\Phi_{ij}$, the dissipation tensor and the turbulent transport $D_{ij}$
are not. The closure has moved one level up the hierarchy, exactly as with
Grad's equations relative to Navier--Stokes--Fourier. Algebraic stress
models sit between the two levels: they expand the anisotropy in a tensor
basis built from mean strain and rotation, with scalar coefficients that
depend on invariants. That representation is the one tensor-basis neural
networks later adopted (Section~\ref{sec:learned}).

\subsection{Large-eddy simulation}

LES resolves the large eddies and models the subgrid stress \eqref{eq:sgs}.
Two families of model reflect two goals.

\emph{Functional models} aim at the right net energy transfer to the
unresolved scales. The Smagorinsky model \cite{Smagorinsky63} is
\begin{equation}
  \tau_{ij}-\tfrac13\tau_{kk}\delta_{ij}=-2\nu_t\bar S_{ij},
  \qquad
  \nu_t=(C_s\Delta)^2|\bar S|,
  \qquad |\bar S|=\sqrt{2\bar S_{ij}\bar S_{ij}} .
  \label{eq:smag}
\end{equation}
Its transfer $-\tau_{ij}\bar S_{ij}=\nu_t|\bar S|^2\ge0$ always removes
energy from the resolved field. Requiring that it remove energy at the
rate of the Kolmogorov cascade gives
$C_s=\pi^{-1}(2/3C_K)^{3/4}\approx0.17$ for $C_K=1.5$
(Exercise~\ref{ex:smag}). In shear flows smaller values are needed, which
is the motivation for the dynamic procedure.

\emph{The dynamic procedure} \cite{Germano91,Lilly92} computes the
coefficient from the resolved field. Apply a second, wider test filter
(hat) and let $T_{ij}$ be the subgrid stress at the test level. Then
\begin{equation}
  L_{ij}\equiv T_{ij}-\widehat{\tau_{ij}}
  =\widehat{\bar u_i\bar u_j}-\hat{\bar u}_i\hat{\bar u}_j
\end{equation}
is computable from the resolved velocity (the Germano identity).
Inserting \eqref{eq:smag} at both levels with the same coefficient
$C=C_s^2$ gives $L_{ij}^{d}=C\,M_{ij}$ with
$M_{ij}=2\Delta^2\widehat{|\bar S|\bar S_{ij}}
-2\hat\Delta^2|\hat{\bar S}|\hat{\bar S}_{ij}$, and the least-squares
solution is $C=\avg{L_{ij}M_{ij}}/\avg{M_{ij}M_{ij}}$. The idea is general:
the closure is required to be consistent across two levels of
coarse-graining, and its free parameter is fixed by that requirement.

\emph{Structural models} aim at the stress tensor itself. The gradient
model $\tau_{ij}\approx\tfrac1{12}\Delta^2\,\partial_k\bar u_i\,\partial_k\bar u_j$
is the leading term of a Taylor expansion of \eqref{eq:sgs}; scale
similarity and approximate deconvolution are related. In \emph{implicit}
LES no stress is written and the dissipation of the numerical scheme is the
closure, so the discretization is part of the model.

\subsection{The lesson of a priori tests}

Clark, Ferziger and Reynolds \cite{Clark79} filtered direct simulations and
compared modelled and exact subgrid stresses point by point. The
Smagorinsky stress correlates poorly with the exact one, and structural
models correlate well. Embedded in a simulation the ranking reverses: the
Smagorinsky model gives stable runs with reasonable statistics, and purely
structural models tend to be unstable because they do not dissipate enough.
The embedded behaviour depends on the net energy transfer, which a
pointwise correlation does not measure. This is the origin of the rule that
a closure must be tested both ways, and it is inequality \eqref{eq:gronwall}
read in the other direction.

\subsection{Statistical closure theories}
\label{sec:statclosure}

For homogeneous turbulence one can write the hierarchy for the spectrum
$E(k,t)$ directly: its equation involves triple correlations, whose
equation involves fourth-order ones.
\begin{itemize}
  \item \textbf{Quasi-normal.} Discard the fourth cumulant, as in
    Example~\ref{ex:langevin}. The resulting equations develop negative
    energy spectra: the closure is not realizable.
  \item \textbf{EDQNM} \cite{Orszag70}. Add a damping of the triple
    correlations at an eddy-turnover rate and make the equation Markovian.
    Realizable, with a free constant.
  \item \textbf{Direct-interaction approximation}
    \cite{Kraichnan59}. A self-consistent closure for the two-time
    correlation and the response function, with no free constant. It is the
    exact statistics of a model system and therefore realizable. In its
    Eulerian form it does not reproduce the Kolmogorov spectrum, because it
    does not separate sweeping of small eddies by large ones from their
    distortion; Lagrangian versions repair this.
\end{itemize}
The response function introduced by the direct-interaction approximation is
the same object as the memory kernel of Section~\ref{sec:memory}, and the
plasma versions (resonance broadening, Section~\ref{sec:waves}) are
closures of this type.

\section{Memory closures}
\label{sec:memory}

The Mori--Zwanzig formalism \cite{Mori65,Zwanzig61,Zwanzig} projects the
full dynamics onto the retained variables without approximation. It does
not produce a closure. It shows what every closure approximates.

\subsection{The linear case is a Schur complement}

Split a linear system into retained variables $x$ and discarded variables
$y$,
\begin{equation}
  \frac{\dd}{\dd t}\begin{pmatrix}x\\y\end{pmatrix}
  =\begin{pmatrix}A&B\\C&D\end{pmatrix}\begin{pmatrix}x\\y\end{pmatrix}.
\end{equation}
Solving the second row for $y$ and substituting gives
\begin{equation}
  \dot x(t)=A\,x(t)+\int_0^tB\,e^{Ds}C\,x(t-s)\,\dd s+B\,e^{Dt}y(0),
  \label{eq:mz-linear}
\end{equation}
which is \eqref{eq:fastslow-exact} with matrices. Under the Laplace
transform $\hat x(s)=\int_0^\infty e^{-st}x(t)\,\dd t$,
\begin{equation}
  \bigl[s-A-B(s-D)^{-1}C\bigr]\hat x(s)=x(0)+B(s-D)^{-1}y(0) .
  \label{eq:schur}
\end{equation}
The operator on the left is the Schur complement of the resolvent. The
exact closure is the frequency-dependent matrix
$\hat K(s)=B(s-D)^{-1}C$.

Every closure of a linear problem is an approximation of $\hat K(s)$:
\begin{itemize}
  \item \textbf{Quasi-static (Markov):} $\hat K(s)\approx\hat K(0)=-BD^{-1}C$.
    This is the Chapman--Enskog closure. It is accurate when the retained
    dynamics is slow compared with the spectrum of $D$.
  \item \textbf{First correction:}
    $\hat K(s)\approx-BD^{-1}C-s\,BD^{-2}C$, which gives
    $(I+BD^{-2}C)\dot x=(A-BD^{-1}C)x$, the rational closure of
    Example~\ref{ex:fastslow}.
  \item \textbf{Rational in $s$:} approximate $\hat K$ by a ratio of
    polynomials. This is an extended-state model
    (Section~\ref{sec:embedding}).
\end{itemize}

\subsection{The general identity}

For a nonlinear system $\dot X=F(X)$ let $\cL=F\cdot\nabla$ be the
Liouville operator acting on observables, $\cP$ a projection onto functions
of the retained variables, and $\cQ=1-\cP$. The Dyson identity
\[
  e^{t\cL}=e^{t\cQ\cL}+\int_0^te^{(t-s)\cL}\,\cP\cL\,e^{s\cQ\cL}\,\dd s
\]
applied to $\cQ\cL a$ for a retained observable $a$ gives
\begin{equation}
  \frac{\dd}{\dd t}a(t)=e^{t\cL}\cP\cL a
    +\int_0^te^{(t-s)\cL}\,\cP\cL\,e^{s\cQ\cL}\cQ\cL a\,\dd s
    +e^{t\cQ\cL}\cQ\cL a .
  \label{eq:mz}
\end{equation}
The three terms are a Markovian term, a memory term, and the
orthogonal dynamics $F(t)=e^{t\cQ\cL}\cQ\cL a$, which depends on the
discarded initial data and satisfies $\cP F(t)=0$. With Mori's linear
projection onto $a$ in an equilibrium inner product, \eqref{eq:mz} becomes
the generalized Langevin equation
\begin{equation}
  \frac{\dd a}{\dd t}=i\Omega\,a(t)-\int_0^tK(s)\,a(t-s)\,\dd s+F(t),
  \qquad
  K(t)=\avg{F(t),F(0)}\avg{a,a}^{-1} .
  \label{eq:gle}
\end{equation}
The second relation is the fluctuation--dissipation theorem: memory and
noise are two views of the same discarded dynamics, so a closure that keeps
dissipation and drops the noise is correct only for averages.

The identity is exact and not useful as it stands, because the orthogonal
dynamics $e^{t\cQ\cL}$ is as hard as the original problem. What it provides
is a classification.

\subsection{Approximations of the memory term}

\begin{itemize}
  \item \textbf{Short memory.} $K(t)\approx2\gamma\,\delta(t)$. Markovian
    closure. Valid under time-scale separation.
  \item \textbf{Long memory} ($t$-model \cite{ChorinHaldKupferman02}). Take
    the integrand constant in $s$, so the memory term is $t$ times its
    initial value. Suited to problems with no separation.
  \item \textbf{Finite memory.} Truncate the integral at a fixed lag, or
    expand the kernel in a finite basis. The length of the lag is a
    modelling choice that can be tested (hypothesis 3 of the mathematical
    closure program).
  \item \textbf{Data-driven kernels.} Fit $K$ from time series of the
    retained variables \cite{Lin21}.
\end{itemize}

\subsection{Extended states are memory with exponential kernels}
\label{sec:embedding}

If the kernel is a sum of exponentials, $K(t)=\sum_jc_je^{-\lambda_jt}$,
each term can be generated by one auxiliary variable
$\dot z_j=-\lambda_jz_j+c_j\,a$, and the non-Markovian equation for $a$
becomes a Markovian system for $(a,z_1,\dots)$. This is called Markovian
embedding. It runs both ways.

\begin{example}[Cattaneo's law]
\label{ex:cattaneo}
Retaining the heat flux with a relaxation time $\tau$,
\begin{equation}
  \tau\,\partial_tq+q=-\kappa\,\partial_xT
  \qquad\Longleftrightarrow\qquad
  q(t)=-\frac{\kappa}{\tau}\int_0^\infty e^{-s/\tau}\,\partial_xT(t-s)\,\dd s .
\end{equation}
With time dependence $e^{-i\omega t}$ the conductivity is
$\kappa/(1-i\omega\tau)$. Fourier's law is the limit $\omega\tau\to0$. The
temperature equation changes from parabolic to hyperbolic, with signal
speed $\sqrt{\kappa/\tau}$ in suitable units.
\end{example}

Grad's thirteen-moment system is this construction for the Boltzmann
equation: the stress and heat flux are auxiliary variables whose
elimination produces exponential memory kernels in the conservation laws.
Conversely, whether to keep a flux as a state variable is the question of
whether the memory kernel varies over the frequencies the reduced model
excites. The heat-flux manuscript makes this quantitative
(Section~\ref{sec:repo}): in one of its two embedded tests the exact
quasi-static closure has transient errors of up to $34\%$, and in the other
it stays within $2\%$, and the difference is the time scale of the resolved
dynamics.

\subsection{Continued fractions and the moment hierarchy}
\label{sec:cf}

Applying the projection repeatedly, first onto $a$, then onto the part of
$\dot a$ orthogonal to $a$, and so on, generates an orthogonal basis in
which the dynamics is tridiagonal, and expresses the Laplace transform of
the correlation function as a continued fraction (Mori). The construction
is the Lanczos algorithm.

For free streaming, \eqref{eq:linkin} with $\Lambda=0$, this basis is the
Hermite basis and the tridiagonal system is \eqref{eq:hermitehier}. For a
single Fourier mode define $\tau_n(s)$ by
$\hat a_{n}=-ik\sqrt n\,\hat a_{n-1}/\tau_n(s)$ for $n\ge2$. Then
\begin{equation}
  \tau_n(s)=s+\frac{(n+1)\,k^2}{\tau_{n+1}(s)},
  \label{eq:cf}
\end{equation}
and unrolling it gives the classical continued fraction for the Laplace
transform of the free-streaming correlation $e^{-k^2t^2/2}$,
\[
  \int_0^\infty e^{-st}e^{-k^2t^2/2}\,\dd t
  =\cfrac{1}{s+\cfrac{k^2}{s+\cfrac{2k^2}{s+\cfrac{3k^2}{s+\cdots}}}} .
\]
A closure of the moment hierarchy at level $N$ is a choice of
\emph{terminator} for this fraction: a replacement for the exact tail
$\tau_{N+1}(s)$. Truncation takes $\tau_{N+1}=\infty$. Relaxation closures
take $\tau_{N+1}=s+\nu$. The Hammett--Perkins closures take
$\tau_{N+1}$ equal to its exact value at zero frequency, which is
proportional to $|k|$ (Section~\ref{sec:collisionless}). The moment
hierarchy of kinetic theory and the Mori--Zwanzig hierarchy are the same
object.

\section{Collisional plasmas}
\label{sec:plasma}

Each of the three facts of Section~\ref{sec:easy} fails in a plasma. This
section covers the particle-to-kinetic closure and the collisional
kinetic-to-fluid closure. The collisionless and magnetized cases follow in
Sections~\ref{sec:collisionless} and~\ref{sec:mag}. A standard reference
for the collisional theory is Helander and Sigmar \cite{HelanderSigmar02}.

\subsection{Parameters}

For electrons of density $n$ and temperature $T$ (energy units),
\[
  \lambda_D=\sqrt{\frac{\varepsilon_0T}{ne^2}},\qquad
  \omega_p=\sqrt{\frac{ne^2}{\varepsilon_0m_e}},\qquad
  \Lambda=\frac{4\pi}{3}\,n\lambda_D^3,\qquad
  \Gamma=\frac{e^2}{4\pi\varepsilon_0\,a\,T},\quad
  a=\Bigl(\frac{3}{4\pi n}\Bigr)^{1/3}.
\]
$\Lambda$ is the number of particles in a Debye sphere and $\Gamma$ is the
ratio of the potential energy of neighbours to the thermal energy. They are
one parameter: $\Gamma=\tfrac13\Lambda^{-2/3}$
(Exercise~\ref{ex:gamma}). A weakly coupled plasma has $\Lambda\gg1$ and
$\Gamma\ll1$.

\subsection{Particles to kinetic: mean field first, collisions second}

The Coulomb force is long-range, so each particle interacts with many
others at once. When $\Lambda\gg1$ the leading-order closure of BBGKY is a
mean field and not a collision term:
\begin{equation}
  \partial_t f_s + \bv\cdot\nabla f_s
  + \frac{q_s}{m_s}\left(\bE + \bv\times\bB\right)\cdot\nabla_{\bv} f_s
  = \sum_{s'}C_{ss'}[f_s,f_{s'}],
  \label{eq:vlasov-landau}
\end{equation}
with $\bE$ and $\bB$ determined self-consistently by Maxwell's equations.
Collisions enter as a correction of relative size $\ln\Lambda/\Lambda$.
They are the cumulative effect of many small-angle deflections, so the
operator is of Fokker--Planck type. For like particles the Landau operator
\cite{Landau36} is
\begin{equation}
  C[f,f]=\frac{q^4\ln\Lambda}{8\pi\varepsilon_0^2m^2}\;
  \nabla_{\bv}\cdot\!\int\mathbf U(\bv-\bv')\cdot
  \bigl[f(\bv')\nabla_{\bv}f(\bv)-f(\bv)\nabla_{\bv'}f(\bv')\bigr]\dd^3v',
  \qquad
  \mathbf U(\bw)=\frac{|\bw|^2\mathbf I-\bw\bw}{|\bw|^3}.
\end{equation}
The logarithm comes from cutting off the Coulomb interaction at the Debye
length and at the distance of closest approach. The Lenard--Balescu
operator \cite{Lenard60,Balescu60} removes the long-range cutoff by
including the dielectric response of the plasma, and is obtained from BBGKY
by neglecting $g_3$ and solving the equation for $g_2$. The expansion
parameter for the particle-to-kinetic closure is weak coupling, not low
density. A rigorous derivation of these equations from particle dynamics is
open; the strategy explored in \texttt{hani-to-algorithms/} is one
approach.

When $\Gamma\gtrsim1$, as for the dust component of a dusty plasma, pair
correlations are not small and this closure fails
(Section~\ref{sec:strong}).

\subsection{Collisions weaken with energy}
\label{sec:coulomb-tail}

The Coulomb collision frequency falls steeply with speed. For electrons
scattering on ions of charge $Z$,
\begin{equation}
  \nu_{ei}(v)=\frac{n_iZ^2e^4\ln\Lambda}{4\pi\varepsilon_0^2m_e^2\,v^3},
  \qquad
  \lambda(v)=\frac{v}{\nu(v)}\propto v^4 .
\end{equation}
There is no single mean free path. The Chapman--Enskog calculation still
goes through (Spitzer and H\"arm \cite{SH53}; Braginskii
\cite{Braginskii65}) and gives $\bq=-\kappa\nabla T$ with
$\kappa\propto T^{5/2}/(Z\ln\Lambda)$ times a function of $Z$. But the
integrand of the heat flux weights high speeds heavily: the heat is carried
by electrons at several times the thermal speed, whose mean free path is
two orders of magnitude longer than the thermal one. The local closure
therefore fails when the thermal mean free path is still of order one
percent of the temperature gradient length or less. In the linear response
language of \eqref{eq:deficit}, the coefficient $K$ is large.

The one-dimensional model results in this repository isolate the mechanism
(Section~\ref{sec:between}). With a velocity-independent collision
frequency the local closure fails only at $k\lambda$ of order one. A
$|v|^{-3}$ tail in the collision frequency moves the breakdown to
$k\lambda\sim10^{-2}$.

\subsection{Repairs used in practice}

\begin{itemize}
  \item \textbf{Flux limiter.} Cap the local flux at a fraction of the
    free-streaming flux,
    \begin{equation}
      q=\min\bigl(q_{\mathrm{SH}},\;f\,nT\vt\bigr)
      \quad\text{or}\quad
      \frac1q=\frac1{q_{\mathrm{SH}}}+\frac1{f\,nT\vt},
    \end{equation}
    with $f$ tuned between about $0.03$ and $0.15$ against experiments or
    kinetic simulations. It is local, it has no predictive value for $f$,
    and it cannot produce preheat ahead of a temperature front.
  \item \textbf{Convolution kernels.} Luciani, Mora and Virmont
    \cite{LMV83} write
    \begin{equation}
      q(x)=\int W(x,x')\,q_{\mathrm{SH}}(x')\,\dd x',
    \end{equation}
    with an exponential kernel whose range is a few tens of thermal mean
    free paths. In a uniform background an exponential kernel of range
    $\ell$ is the multiplier $1/(1+k^2\ell^2)$, which has the form
    \eqref{eq:rational}.
  \item \textbf{Fits to linear kinetic theory.} Epperlein and Short
    \cite{ES91} built a nonlocal model on linear Fokker--Planck results.
    The fit usually quoted from that line of work is
    $\kappa/\kappa_{\mathrm{SH}}=1/\bigl(1+(30\,k\lambda_e)^{4/3}\bigr)$,
    with $\lambda_e=\sqrt{\lambda_{ee}\lambda_{ei}}$ a delocalization length
    and not the thermal mean free path. The formula was confirmed from a
    later paper that quotes it, not from the original.
  \item \textbf{Multigroup diffusion.} Schurtz, Nicola\"i and Busquet
    \cite{SNB00} solve one screened diffusion equation per electron energy
    group and sum the contributions. This is the form used in
    radiation-hydrodynamics codes, because it works in several dimensions
    and avoids an explicit convolution.
  \item \textbf{Moment-based integral closures} for heat flux, viscosity
    and friction together at arbitrary collisionality \cite{JiHeld14}.
\end{itemize}
Comparisons against kinetic solvers show that the ranking of these models
depends on the drive, the temperature contrast and the collision physics
\cite{Brodrick17}. The heat-flux manuscript in this repository proves that
for a thermal drive the exact linear closure is a positive superposition of
screened diffusion responses (a Stieltjes function of $k^2$), which is the
structure the multigroup models assume; it does not tell how many groups
are needed or what their weights are (Section~\ref{sec:repo}).

\subsection{Two fluids: more to close}

A plasma has at least two species with a mass ratio of 1836 or more, which
can have different temperatures and flows, and electromagnetic fields that
are dynamical variables. The fluid equations for each species contain, in
addition to $\boldsymbol\pi_s$ and $\bq_s$, the collisional exchange of
momentum $\mathbf R$ and of heat between species. Braginskii's closure
\cite{Braginskii65} for the unmagnetized direction gives, for $Z=1$,
\begin{equation}
  \mathbf R_\parallel=-0.51\,\frac{m_en}{\tau_e}\,u_\parallel
      -0.71\,n\nabla_\parallel T_e,
  \qquad
  q_{e\parallel}=0.71\,nT_e\,u_\parallel
      -3.16\,\frac{nT_e\tau_e}{m_e}\nabla_\parallel T_e ,
\end{equation}
where $u_\parallel$ is the electron flow relative to the ions and $\tau_e$
the electron collision time. The thermal force (second term of
$\mathbf R$) and the convective heat flux (first term of $q$) carry the
same coefficient $0.71$. That is an Onsager symmetry, a structural
constraint that any closure of the coupled system should respect. The
electron momentum equation with $\mathbf R$ closed is the generalized Ohm's
law (Section~\ref{sec:mag}).

The Debye length, the skin depths and the gyroradii of each species are
distinct scales, and a closure valid at one may fail at another.

\subsection{Closures in sequence}

A plasma model stacks several closures, each with its own defect:
\begin{center}
\renewcommand{\arraystretch}{1.3}
\begin{tabularx}{\textwidth}{@{}P{4.2cm}P{4.2cm}Y@{}}
\toprule
Transition & Small parameter & Fails when \\
\midrule
Particles $\to$ kinetic & $1/\Lambda$ (weak coupling)
   & strong coupling, dense or dusty plasmas \\
Kinetic $\to$ fluid, collisional & $\lambda/L$
   & steep gradients, fast particles \\
Kinetic $\to$ fluid, magnetized & $\rho/L$
   & unmagnetized regions, parallel dynamics \\
Kinetic $\to$ fluid, collisionless & none; matched to linear response
   & nonlinear trapping, strong drives \\
Fluid $\to$ mean transport & none in general
   & always a model (turbulence) \\
\bottomrule
\end{tabularx}
\end{center}

\section{Collisionless closures}
\label{sec:collisionless}

In hot, dilute plasmas the mean free path exceeds the system size. Nothing
relaxes $f$ toward a Maxwellian, and Chapman--Enskog has no small
parameter. What takes the place of collisions is \emph{phase mixing}:
particles streaming at different speeds smear out perturbations, and
particles resonant with a wave ($v=\omega/k$) exchange energy with it
(Landau damping \cite{Landau46}). This section derives the closure that
builds phase mixing into a fluid model, following Hammett and Perkins
\cite{HP90,HDP92}. The derivation is the one in the companion note, with
the same conventions; the closure at general order
(Section~\ref{sec:general-order}) and the error figures are new here.

\subsection{Setup and conventions}

One species, one spatial dimension $z$ (along the magnetic field), no
collisions, linearized about
\[
  f_0(v)=\frac{n_0}{\sqrt{2\pi}\,\vt}\exp\Bigl(-\frac{v^2}{2\vt^2}\Bigr),
  \qquad \vt=\sqrt{T_0/m}.
\]
With acceleration $a=(e/m)E$ the perturbation obeys
\begin{equation}
  \partial_tf_1+v\,\partial_zf_1+a\,\partial_vf_0=0 .
  \label{eq:vlasov}
\end{equation}
Derivations use units $n_0=m=\vt=1$, in which this is \eqref{eq:linkin}
with $\Lambda=0$; dimensional forms are restored for results. Fourier modes
are $\propto e^{ikz-i\omega t}$, and
\begin{equation}
  \zeta=\frac{\omega}{\sqrt2\,|k|\,\vt}.
  \label{eq:zeta}
\end{equation}
The Hilbert transform is taken with the convention
\begin{equation}
  \Hil[g](z)=\frac1\pi\,\PV\int_{-\infty}^{\infty}\frac{g(z')}{z-z'}\,\dd z',
  \qquad
  \widehat{\Hil[g]}(k)=-i\,\sgn(k)\,\hat g(k).
  \label{eq:hilbert}
\end{equation}

\subsection{The kinetic response and the origin of $\sgn(k)$}

In Fourier space \eqref{eq:vlasov} gives, using $\partial_vf_0=-vf_0$,
$f_{1,k}=-i\,a\,v\,f_0/(kv-\omega)$, so the density perturbation is
\[
  n_k=-\frac{ia}{k}\int\frac{v\,f_0}{v-\omega/k}\,\dd v
     =-\frac{ia}{k}\Bigl[1+\frac\omega k\int\frac{f_0}{v-\omega/k}\,\dd v\Bigr].
\]
The integral is defined for $\operatorname{Im}\omega>0$ (causality) and
continued analytically. It is the plasma dispersion function
\begin{equation}
  Z(\zeta)=\frac1{\sqrt\pi}\int_{-\infty}^{\infty}\frac{e^{-x^2}}{x-\zeta}\,\dd x,
  \qquad\operatorname{Im}\zeta>0,
  \label{eq:Z}
\end{equation}
the Hilbert transform, in velocity space, of a Gaussian. For real $\zeta$
the Plemelj formula
$1/(x-\zeta-i0)=\PV\,1/(x-\zeta)+i\pi\,\delta(x-\zeta)$ separates a
principal-value part from the contribution of the resonant particles.

The pole is at $v=\omega/k$. For $\operatorname{Im}\omega>0$ it lies above
the real $v$ axis when $k>0$ and below when $k<0$. Substituting $x\to-x$ in
the second case shows that both signs give the same function of $\zeta$ as
defined in \eqref{eq:zeta}, with $|k|$ and not $k$:
\begin{equation}
  n_k=\frac{a}{ik}\,R(\zeta),\qquad R(\zeta)=1+\zeta Z(\zeta).
  \label{eq:R}
\end{equation}
This is the source of the non-analytic dependence on $k$. The expansions
needed are
\begin{align}
  R(\zeta)&=1+i\sqrt\pi\,\zeta-2\zeta^2+O(\zeta^3), &\zeta&\to0,
     \label{eq:Rsmall}\\
  R(\zeta)&=-\frac1{2\zeta^2}-\frac3{4\zeta^4}-\frac{15}{8\zeta^6}-\dots,
     &|\zeta|&\to\infty. \label{eq:Rlarge}
\end{align}
The term $i\sqrt\pi\,\zeta$ is the resonant contribution,
$Z(0)=i\sqrt\pi$. On the real axis $|R|\le1$.

\subsection{Linearized moment equations}

Define $n=\int f_1\,\dd v$, $u=\int vf_1\,\dd v$, $p=\int v^2f_1\,\dd v$,
$T=p-n$, the heat flux $q=\int v^3f_1\,\dd v-3u$, and the deviation of the
fourth moment from its Maxwellian value,
$\delta r=\int v^4f_1\,\dd v-6p+3n$. In the Hermite variables of
Section~\ref{sec:hermite}, $(n,u,T,q,\delta r)
=(a_0,a_1,\sqrt2a_2,\sqrt6a_3,\sqrt{24}a_4)$. Moments of \eqref{eq:vlasov}
give
\begin{align}
  \partial_tn+\partial_zu&=0, \label{eq:n}\\
  \partial_tu+\partial_z(n+T)&=a, \label{eq:u}\\
  \partial_tT+2\,\partial_zu+\partial_zq&=0, \label{eq:T}\\
  \partial_tq+3\,\partial_zT+\partial_z\delta r&=0. \label{eq:q}
\end{align}
Each is exact. In Fourier space write
$\omega/k=\sqrt2\,\sgn(k)\,\zeta$.

\subsection{Three-moment closure}

Keep \eqref{eq:n}--\eqref{eq:T} and close $q$. A diffusive closure is
$q_k=-\chi\,ik\,T_k$. With collisions $\chi\sim\vt^2/\nu$. Without them the
only available rate is the phase-mixing rate $|k|\vt$, so
$\chi=\chi_1\sqrt2\,\vt/|k|$ with $\chi_1$ a pure number:
\begin{equation}
  q_k=-i\,\sgn(k)\,\sqrt2\,\chi_1\,T_k .
  \label{eq:closure3}
\end{equation}
Equation \eqref{eq:T} then gives $T_k=2\zeta\,n_k/(\zeta+i\chi_1)$, and
\eqref{eq:u} gives $n_k=(a/ik)\,R_3$ with
\begin{equation}
  R_3(\zeta)=\frac{\chi_1-i\zeta}
     {\chi_1-3i\zeta-2\chi_1\zeta^2+2i\zeta^3}.
  \label{eq:R3}
\end{equation}
For large $\zeta$, $R_3=-1/(2\zeta^2)-3/(4\zeta^4)+O(\zeta^{-5})$, which
agrees with \eqref{eq:Rlarge} for any $\chi_1$ because the retained moment
equations are exact. For small $\zeta$,
$R_3=1+2i\zeta/\chi_1+O(\zeta^2)$, and comparison with \eqref{eq:Rsmall}
fixes
\begin{equation}
  \chi_1=\frac2{\sqrt\pi}.
\end{equation}
Restoring units,
\begin{equation}
  q_k=-n_0\,\chi_1\frac{\sqrt2\,\vt}{|k|}\,ik\,T_k
     =-i\,\sgn(k)\;n_0\sqrt{\frac8\pi}\;\vt\;T_k .
  \label{eq:hp3}
\end{equation}
By \eqref{eq:hilbert} the multiplier $-i\,\sgn(k)$ is the Hilbert
transform, so in real space
\begin{equation}
  q(z)=n_0\sqrt{\frac8\pi}\;\vt\;\Hil[T](z)
      =-n_0\,\chi_1\frac{\sqrt2\,\vt}{\pi}
        \int_0^\infty\frac{T(z+z')-T(z-z')}{z'}\,\dd z' .
  \label{eq:real3}
\end{equation}
The kernel carries no length scale, because a collisionless plasma has no
mean free path to supply one.

Compare the two heat-flux closures:
\begin{equation}
  \text{collisional:}\quad q_k=-\kappa\,ik\,T_k,
  \qquad\qquad
  \text{collisionless:}\quad
  q_k=-n_0\sqrt{\frac8\pi}\;\vt\;i\,\sgn(k)\,T_k .
\end{equation}
The first damps temperature perturbations at a rate $\propto k^2$, the
second at a rate $\propto|k|\vt$. In real space the first is a derivative
and the second a nonlocal operator. Both are dissipative: with
\eqref{eq:hp3},
$\int q\,\partial_zT\,\dd z=-n_0\sqrt{8/\pi}\,\vt\sum_k|k|\,|T_k|^2\le0$.

\subsection{Four-moment closure}

Now keep \eqref{eq:q} as well and close $\delta r$ in terms of the two
highest retained moments. Parity and dimensions allow a local term in $T$
and a phase-mixing term in $q$:
\begin{equation}
  \delta r_k=\beta_1\,T_k-i\,\sgn(k)\,\sqrt2\,D_1\,q_k .
  \label{eq:closure4}
\end{equation}
The $T$ term is not dissipative; it replaces the coefficient $3$ in
\eqref{eq:q} by $b\equiv3+\beta_1$. Eliminating upward through the
hierarchy,
\begin{align*}
  \text{from \eqref{eq:q}:}\quad
    q_k&=\frac{b\,T_k}{\sqrt2\,\sgn(k)\,(\zeta+iD_1)},\\[4pt]
  \text{from \eqref{eq:T}:}\quad
    T_k&=\frac{4\zeta(\zeta+iD_1)}{2\zeta(\zeta+iD_1)-b}\;n_k,\\[4pt]
  \text{from \eqref{eq:u}:}\quad
    n_k&=\frac{a}{ik}\,R_4(\zeta),
\end{align*}
with
\begin{equation}
  R_4(\zeta)=\frac{b-2iD_1\zeta-2\zeta^2}
     {b-6iD_1\zeta-(2b+6)\zeta^2+4iD_1\zeta^3+4\zeta^4}.
  \label{eq:R4}
\end{equation}
Every factor of $\sgn(k)$ cancels except through $\zeta$.

\paragraph{Matching at small $\zeta$.} With $\alpha=2iD_1/b$,
\[
  R_4(\zeta)=1+2\alpha\,\zeta
    +\Bigl(2+\frac4b+6\alpha^2\Bigr)\zeta^2+O(\zeta^3).
\]
Comparing with \eqref{eq:Rsmall}:
\begin{align*}
  O(\zeta):&\quad2\alpha=i\sqrt\pi
    \;\Longrightarrow\;D_1=\frac{\sqrt\pi}{4}\,b,\\
  O(\zeta^2):&\quad2+\frac4b-\frac{3\pi}{2}=-2
    \;\Longrightarrow\;b=\frac{8}{3\pi-8}.
\end{align*}
Therefore
\begin{equation}
  D_1=\frac{2\sqrt\pi}{3\pi-8}\approx2.49,
  \qquad
  \beta_1=\frac{32-9\pi}{3\pi-8}\approx2.62 .
  \label{eq:coeffs4}
\end{equation}

\paragraph{Behaviour at large $\zeta$.} From \eqref{eq:R4},
\[
  \frac1{R_4}=-2\zeta^2+3+\frac{b}{\zeta^2}+O(\zeta^{-3}),
  \qquad\text{whereas}\qquad
  \frac1R=-2\zeta^2+3+\frac3{\zeta^2}+\dots
\]
So $R_4$ reproduces the first two terms of \eqref{eq:Rlarge} for any $D_1$
and $\beta_1$, but the $\zeta^{-6}$ term is correct only if $\beta_1=0$.
The closure spends its two parameters on the low-frequency side and gives
up one order of high-frequency accuracy to do so.

\paragraph{Real space.} Restoring units,
\begin{equation}
  \delta r(z)=\beta_1\,n_0\vt^2\,T(z)+\sqrt2\,D_1\,\vt\,\Hil[q](z),
\end{equation}
to be used in
$\partial_tq+3n_0\vt^2\,\partial_zT+\partial_z\delta r=0$. The Hilbert
transform now acts on the heat flux, which has its own evolution equation.

\subsection{How accurate the closures are}

\begin{center}
\renewcommand{\arraystretch}{1.5}
\begin{tabular}{lll}
\toprule
  & Three-moment & Four-moment \\
\midrule
  Closed quantity & $q$ & $\delta r$ \\
  Closure & $q=n_0\sqrt2\,\chi_1\vt\,\Hil[T]$
          & $\delta r=\beta_1n_0\vt^2T+\sqrt2\,D_1\vt\,\Hil[q]$ \\
  Coefficients & $\chi_1=2/\sqrt\pi$
          & $D_1=\dfrac{2\sqrt\pi}{3\pi-8}$,\;
            $\beta_1=\dfrac{32-9\pi}{3\pi-8}$ \\
  Matched as $\zeta\to0$ & through $O(\zeta)$ & through $O(\zeta^2)$ \\
  Matched as $\zeta\to\infty$ & through $O(\zeta^{-4})$
          & through $O(\zeta^{-4})$ \\
  $\sup_{\zeta\in\mathbb R}|R_n-R|$ & $0.34$ at $\zeta\approx1.0$
          & $0.059$ at $\zeta\approx1.3$ \\
\bottomrule
\end{tabular}
\end{center}
The last row was computed for this note; Figure~\ref{fig:response} shows
the functions. Both closures are exact at the two ends of the frequency
axis and are rational (Pad\'e-type) approximants in between. The
three-moment closure is poor at phase velocities near the thermal speed,
where most particles are resonant; the four-moment closure is within six
percent of the peak response everywhere on the real axis.

\begin{figure}[ht]
\centering
\begin{tikzpicture}
\begin{axis}[width=0.48\textwidth,height=5.6cm,xlabel={$\zeta$},
  title={$\operatorname{Re}R$},xmin=0,xmax=4,grid=major,
  legend style={font=\footnotesize,at={(0.97,0.97)},anchor=north east}]
\addplot[black,thick] table[x=zeta,y=ReR]{\rundir/response.dat};
\addplot[blue,dashed,thick] table[x=zeta,y=ReR3]{\rundir/response.dat};
\addplot[red,densely dotted,very thick] table[x=zeta,y=ReR4]{\rundir/response.dat};
\legend{exact,3-moment,4-moment}
\end{axis}
\end{tikzpicture}\hfill
\begin{tikzpicture}
\begin{axis}[width=0.48\textwidth,height=5.6cm,xlabel={$\zeta$},
  title={$\operatorname{Im}R$},xmin=0,xmax=4,grid=major]
\addplot[black,thick] table[x=zeta,y=ImR]{\rundir/response.dat};
\addplot[blue,dashed,thick] table[x=zeta,y=ImR3]{\rundir/response.dat};
\addplot[red,densely dotted,very thick] table[x=zeta,y=ImR4]{\rundir/response.dat};
\end{axis}
\end{tikzpicture}
\caption{The collisionless density response $R(\zeta)=1+\zeta Z(\zeta)$ on
the real axis and its three- and four-moment approximants \eqref{eq:R3},
\eqref{eq:R4}. Data: \texttt{\rundir/response.dat}.}
\label{fig:response}
\end{figure}

\subsection{The closure at general order}
\label{sec:general-order}

The three-moment closure has a simple description that extends to any
order. In the Hermite hierarchy \eqref{eq:hermitehier} with $\Lambda=0$,
the forced response at frequency $\omega$ has coefficients, for $n\ge1$,
\[
  \hat a_n\;\propto\;\frac{W_n(z)}{\sqrt{n!}},\qquad
  W_n(z)=\Bigl\langle\frac{\He_n(v)}{v-z}\Bigr\rangle,\qquad
  z=\frac{\omega}{k},
\]
and the Hermite recurrence gives
$W_{n+1}=z\,W_n-n\,W_{n-1}+\delta_{n0}$, with
$W_0(0)=i\sqrt{\pi/2}\,\sgn(k)$ and $W_1(0)=1$. At zero frequency the
ratio of consecutive coefficients is therefore an explicit number:
\begin{equation}
  \hat a_{N+1}=-i\,\sgn(k)\,c_N\,\hat a_N,
  \qquad
  c_N=\frac{N!!}{(N-1)!!\,\sqrt{N+1}}\times
  \begin{cases}\sqrt{2/\pi},&N\text{ even},\\[2pt]
               \sqrt{\pi/2},&N\text{ odd}.\end{cases}
  \label{eq:cN}
\end{equation}
Equivalently $c_Nc_{N+1}=\sqrt{(N+1)/(N+2)}$ with
$c_2=4/\sqrt{6\pi}$. In terms of the continued fraction \eqref{eq:cf},
$c_N=\sqrt{N+1}\,|k|/\tau_{N+1}(0)$: closing with \eqref{eq:cN} replaces
the tail of the continued fraction by its value at zero frequency.

\begin{center}
\begin{tabular}{lrrrrrrr}
\toprule
$N$ & 2 & 3 & 4 & 5 & 6 & 7 & 8 \\
\midrule
$c_N$ & 0.9213 & 0.9400 & 0.9515 & 0.9594 & 0.9650 & 0.9693 & 0.9727 \\
\bottomrule
\end{tabular}
\end{center}
(The formula was checked against direct quadrature of $W_n$ at
$z=10^{-4}i$; agreement is to $2\times10^{-4}$, the accuracy of the
quadrature.)

Three observations.
\begin{itemize}
  \item $N=2$ is Hammett--Perkins:
    $q=\sqrt6\,a_3$ and $T=\sqrt2\,a_2$ give
    $q_k=-i\,\sgn(k)\sqrt3\,c_2\,T_k=-i\,\sgn(k)\sqrt{8/\pi}\,T_k$.
  \item $N=3$ gives $\delta r_k=-i\,\sgn(k)\,\tfrac32\sqrt{\pi/2}\;q_k$,
    which is \eqref{eq:closure4} with $\beta_1=0$ and
    $D_1=\tfrac{3}{4}\sqrt\pi$. It matches $R$ through $O(\zeta)$ only. The
    Hammett--Perkins four-moment closure uses the additional term
    $\beta_1T$ to match $O(\zeta^2)$ as well, at the cost described above.
  \item $c_N\to1$ as $N\to\infty$. The limiting closure
    $\hat a_{N+1}=-i\,\sgn(k)\,\hat a_N$ is an outgoing-wave condition:
    free streaming carries perturbations toward high $n$, and the closure
    lets them leave instead of reflecting at the truncation. This is the
    absorber that Section~\ref{sec:moment} said a truncated system lacks.
\end{itemize}

\subsection{Caveats}
\label{sec:hp-caveats}

\begin{itemize}
  \item \textbf{Thermal-speed convention.} This note and Hammett--Perkins
    use $\vt=\sqrt{T_0/m}$. Hunana et al.\ \cite{Hunana19} use
    $v_{\mathrm{th}}=\sqrt{2T_0/m}$, in which the three-moment closure
    reads $q_k=-i\,(2/\sqrt\pi)\,n_0v_{\mathrm{th}}\,\sgn(k)\,T_k$ and the
    local coefficient of the four-moment closure appears a factor of two
    smaller.
  \item \textbf{Heat-flux convention.} Here $q$ is the third central
    moment. The heat-flux manuscript uses the energy flux, which is half of
    it in one dimension, so $\sqrt{8/\pi}$ becomes $\sqrt{2/\pi}$ there
    (Appendix~\ref{app:conv}).
  \item \textbf{Sign of the Hilbert transform.} With \eqref{eq:hilbert} and
    modes $e^{+ikz}$, $i\,\sgn(k)\leftrightarrow-\Hil$. Several
    Landau-fluid papers define the opposite sign. The integral form
    \eqref{eq:real3} is convention-free.
  \item \textbf{Known inaccuracies.} According to \cite{Hunana19}, the
    three-moment closure gives $q\sim\zeta^{-2}$ at large $\zeta$ where the
    kinetic result is $\zeta^{-3}$, and so tends to overestimate Landau
    damping. The closures are local in time.
  \item \textbf{Linearity.} The matching is linear, about a Maxwellian.
    Collisionless phase mixing is reversible: the perturbation is still in
    the distribution function, at fine velocity scales, and a second
    perturbation can bring it back as a plasma echo \cite{Gould67}. A
    closure that damps cannot do that. Nor does it describe particle
    trapping, which stops Landau damping after a bounce time at finite
    amplitude.
  \item \textbf{Domain.} $\Hil$ is defined on the whole line or a periodic
    domain. Bounded or sheath-terminated field lines need a separate choice
    of extension.
\end{itemize}

\subsection{Extensions used in practice}

\begin{itemize}
  \item \textbf{Three dimensions and magnetization.} The parallel and
    perpendicular pressures are evolved separately, with one closure for
    each parallel heat flux \cite{Snyder97}. Gyrofluid models add closures
    for finite-gyroradius averaging and for drift resonances
    \cite{DorlandHammett93}.
  \item \textbf{Collisions.} The collisionless and collisional rates are
    usually combined by adding them, which is the harmonic interpolant of
    Section~\ref{sec:between}. The exact answer for model collision
    operators is known there and can be compared.
  \item \textbf{Real-space evaluation.} Along curved or bounded field lines
    a Fourier transform is unavailable. Since
    \begin{equation}
      \frac1{|k|}=\frac1\pi\int_0^\infty\frac{m^{-1/2}}{m+k^2}\,\dd m,
      \label{eq:lorentz}
    \end{equation}
    the operator $1/|k|$ can be approximated over a band of wavenumbers by
    a short sum of terms $w_j/(m_j+k^2)$, each of which is a tridiagonal
    solve in real space \cite{DJU14}. Equation \eqref{eq:lorentz} says that
    the collisionless closure is the limiting member of the same family of
    screened-diffusion closures as the collisional nonlocal models of
    Section~\ref{sec:plasma}.
  \item \textbf{Relaxation closures in ten-moment models} replace
    $|k|$ by a fixed wavenumber $k_0$, giving a local term that relaxes the
    pressure tensor toward isotropy at a rate $\vt|k_0|$ \cite{Wang15}.
\end{itemize}

\section{Between the limits: closure as a function of wavenumber and frequency}
\label{sec:between}

Sections~\ref{sec:asym} and~\ref{sec:collisionless} fixed the heat flux at
the two ends: Fourier's law when collisions dominate and the
Hammett--Perkins closure when they are absent. Most applications sit in
between. This section states the problem in a form that covers every
closure met so far, and reports what is known exactly.

\subsection{One function}

For the linear problem \eqref{eq:linkin}, a single Fourier mode, and data
that lie in the retained moments, the exact relation between heat flux and
temperature is a multiplier,
\begin{equation}
  \hat q(k,\omega)=-ik\,\chi(k,\omega)\,\hat T(k,\omega) .
  \label{eq:symbol}
\end{equation}
This is the Schur complement \eqref{eq:schur} for the kinetic equation:
$\chi$ is the exact memory kernel, in Fourier--Laplace variables. Every
linear closure is a choice of approximation to $\chi$:
\begin{center}
\renewcommand{\arraystretch}{1.35}
\begin{tabularx}{\textwidth}{@{}P{4.4cm}P{4.6cm}Y@{}}
\toprule
Closure & $\chi(k,\omega)$ & Character \\
\midrule
Fourier (Chapman--Enskog) & $\kF$ & local \\
Gradient expansion, truncated & $\kF(1-Kk^2)$
  & local; negative for $k^2>1/K$ \\
Rational in $k^2$ & $\kF/(1+Kk^2)$
  & exponential kernel in space \\
Cattaneo, Grad & $\kF/(1-i\omega\tau)$
  & exponential kernel in time \\
Fokker--Planck fit (Section~\ref{sec:plasma}) & $\kF/(1+(30k\lambda_e)^{4/3})$ & fit \\
Multigroup & $\sum_jw_j/(m_j+k^2)$, $w_j\ge0$
  & sum of screened diffusions \\
Hammett--Perkins & $\sqrt{8/\pi}\;n_0\vt/|k|$
  & scale-free nonlocal \\
Harmonic interpolant & $\bigl(1/\kF+|k|/(\sqrt{8/\pi}\,n_0\vt)\bigr)^{-1}$
  & both limits by construction \\
Exact, BGK & $i\,\Phi\bigl((\omega+i\nu)/|k|\bigr)/|k|$ & see below; units $n_0=\vt=1$ \\
\bottomrule
\end{tabularx}
\end{center}
A flux limiter is not in the table because it is nonlinear: it depends on
the amplitude of the temperature gradient and has no multiplier.

Physical requirements become conditions on the function $\chi$:
causality is analyticity in the upper half of the $\omega$ plane;
dissipativity is $\operatorname{Re}\chi\ge0$ on the real axes; a real
kernel requires $\chi(-k,-\bar\omega)=\overline{\chi(k,\omega)}$; and the
two limits are $\chi\to\kF$ as $k\to0$ and $|k|\chi\to$ a constant as
$|k|\to\infty$ at $\omega=0$.

The \emph{quasi-static} closure is $\chi(k,0)$. It is the right object when
the retained dynamics is slow compared with the kinetic relaxation at that
wavenumber, and Section~\ref{sec:repo} states when that holds.

\subsection{Exact results in one velocity dimension}

Two model collision operators can be solved in closed form. Both results
are proved in the notes cited in Section~\ref{sec:repo}; the formulas were
re-evaluated here and checked against an independent Hermite--Galerkin
solve of \eqref{eq:hermitehier} with 1200 modes (agreement to
$10^{-14}$ at $\nu/|k|=1$ and $3$).

\paragraph{BGK with constant $\nu$.} Let $W(z)=\avg{1/(v-z)}$,
$R(z)=1+zW(z)$, and
\[
  \Phi(z)=\frac{z\,\bigl(1-3R(z)+z^2R(z)\bigr)}{1-R(z)+z^2R(z)} .
\]
Then every solution satisfies
$\hat q=\sgn(k)\,\Phi\bigl((\omega+i\nu)/|k|\bigr)\hat T$: the collisional
closure is the collisionless one evaluated at the complex phase velocity
$(\omega+i\nu)/|k|$. At zero frequency, with $y=\nu/|k|$,
\begin{equation}
  \hat q=-i\,\sgn(k)\,\theta(y)\,\hat T,\qquad
  \theta(y)=\frac{1+y^2}{y\,(3+y^2)}\;
  \frac{\avg{\He_3^2/(v^2+y^2)}}{\avg{\He_2^2/(v^2+y^2)}} .
  \label{eq:theta-bgk}
\end{equation}

\paragraph{Dougherty.} With $E(s)=[(1-s)e^s]^{1/y^2}$,
\begin{equation}
  \theta_1(y)=\frac2y\;\frac{\int_0^1s\,E(s)\,\dd s}{\int_0^1E(s)\,\dd s}.
  \label{eq:theta-dougherty}
\end{equation}

Both tend to $\sqrt{8/\pi}=1.5958$ as $y\to0$ (computed: $1.5949$ at
$y=0.002$ for BGK and $1.573$ at $y=0.02$ for Dougherty), and both tend to
Fourier's law as $y\to\infty$, with $\kF=3/\nu$ and $1/\nu$ respectively.
The leading deficits \eqref{eq:deficit} are $K=4/\nu^2$ and
$K=1/(3\nu^2)$ (computed at $y=50$: $3.986$ and $0.3333$).

\subsection{Comparison at matched conductivity}

Two collision operators with the same nominal $\nu$ have different
conductivities, so a fair comparison fixes $\kF$. Use the variable
$s=|k|\,\kF/(n_0\vt)$, in which Fourier's law is $\chi/\kF=1$ and the
collisionless limit is $\chi/\kF=\sqrt{8/\pi}/s$.

\begin{figure}[ht]
\centering
\begin{tikzpicture}
\begin{loglogaxis}[width=0.78\textwidth,height=7.2cm,
  xlabel={$s=|k|\,\kF/(n_0\vt)$},ylabel={$\chi/\kF$},
  xmin=0.01,xmax=100,ymin=0.01,ymax=1.5,grid=major,
  legend style={font=\footnotesize,at={(0.03,0.03)},anchor=south west}]
\addplot[black,thick] table[x=s,y=bgk]{\rundir/intermediate.dat};
\addplot[red,densely dotted,very thick] table[x=s,y=dougherty]{\rundir/intermediate.dat};
\addplot[blue,dashed,thick] table[x=s,y=harmonic]{\rundir/intermediate.dat};
\addplot[gray,thin] table[x=s,y=limits]{\rundir/intermediate.dat};
\legend{exact BGK,exact Dougherty,harmonic interpolant,the two limits}
\end{loglogaxis}
\end{tikzpicture}
\caption{Quasi-static heat-flux closure in one velocity dimension for two
collision operators at equal conductivity, with the harmonic interpolant
between Fourier's law and the Hammett--Perkins closure. Data:
\texttt{\rundir/intermediate.dat}.}
\label{fig:between}
\end{figure}

\begin{center}
\begin{tabular}{rrrr}
\toprule
$s$ & exact BGK & exact Dougherty & harmonic interpolant \\
\midrule
0.10 & 0.9956 & 0.9967 & 0.9410 \\
0.32 & 0.9595 & 0.9685 & 0.8346 \\
1.00 & 0.7570 & 0.7844 & 0.6148 \\
3.16 & 0.3884 & 0.4001 & 0.3354 \\
10.0 & 0.1466 & 0.1484 & 0.1376 \\
\bottomrule
\end{tabular}
\end{center}

Three conclusions, the first two computed here and the third quoted.
\begin{itemize}
  \item \textbf{The interpolant is biased low for these operators.} Over
    $0.01\le s\le100$ the ratio of exact to interpolated conductivity lies
    between $1.006$ and $1.23$ for BGK and between $1.006$ and $1.28$ for
    Dougherty. An interpolant that is exact at both ends is off by a
    quarter in the middle.
  \item \textbf{These two operators nearly agree at matched conductivity.}
    Their collision frequencies do not depend on speed, and the two curves
    differ by a few percent.
  \item \textbf{A speed-dependent collision frequency changes the picture.}
    For BGK with $\nu(v)=\nu_0/(1+|v|^3)$, which has the Coulomb tail, the
    note \texttt{bgk\_velocity\_dependent\_closure.tex} reports
    $\chi/\kF=0.57$ at $s\approx0.7$ and $0.13$ at $s\approx7$, far below
    the curves in Figure~\ref{fig:between}; the leading deficit coefficient
    at matched conductivity is larger by a factor of about fifteen to
    twenty; and the error of the harmonic interpolant changes sign (the
    interpolant is too high by up to $48\%$). That is the behaviour
    reported for the Landau operator \cite{Brodrick17}.
\end{itemize}
The summary is the main finding of the heat-flux manuscript: matching the
two limits and the conductivity does not determine the closure. The
coefficient $K$ in \eqref{eq:deficit} is an independent property of the
collision operator, and it is the first quantity a benchmark should test
after the conductivity itself.

\section{Magnetized plasmas}
\label{sec:mag}

A magnetic field confines particles to gyro-orbits of radius $\rho$ across
the field while leaving them free along it. This introduces a second small
parameter, $\rho/L$, that works even without collisions. The closure
problem becomes anisotropic: the perpendicular direction is closed by
gyromotion, and the parallel direction keeps all the difficulties of
Sections~\ref{sec:plasma} and~\ref{sec:collisionless}.

\subsection{Collisionless, across the field: double-adiabatic laws}

When the fields vary slowly compared with the gyration, each particle
conserves its magnetic moment $m v_\perp^2/2B$, and the pressure tensor is
gyrotropic, $\mathbf P=p_\perp(\mathbf I-\bb\bb)+p_\parallel\bb\bb$ with
$\bb=\bB/B$. If the heat fluxes along the field are neglected, the two
pressures obey the double-adiabatic laws of Chew, Goldberger and Low
\cite{CGL56},
\begin{equation}
  \frac{\dd}{\dd t}\!\left(\frac{p_\perp}{nB}\right)=0,
  \qquad
  \frac{\dd}{\dd t}\!\left(\frac{p_\parallel B^2}{n^3}\right)=0 .
\end{equation}
The first is conservation of the magnetic moment; the second is
conservation of the parallel action for a flux tube whose length scales as
$B/n$. This is an algebraic closure with a real small parameter in the
perpendicular direction and none in the parallel direction: setting the
parallel heat fluxes to zero is the truncation of Section~\ref{sec:moment},
and it omits Landau damping. Landau-fluid models restore it by closing the
parallel heat fluxes as in Section~\ref{sec:collisionless}
\cite{Snyder97}.

\subsection{Collisional: Braginskii}

With collisions and a magnetic field the heat flux has three parts,
\begin{equation}
  \bq=-\kappa_\parallel\nabla_\parallel T-\kappa_\perp\nabla_\perp T
      \mp\kappa_\wedge\,\bb\times\nabla T ,
\end{equation}
and for electrons with $Z=1$ and $\Omega_e\tau_e\gg1$ \cite{Braginskii65}
\[
  \kappa_\parallel=3.16\,\frac{nT_e\tau_e}{m_e},\qquad
  \kappa_\perp=4.66\,\frac{nT_e}{m_e\Omega_e^2\tau_e},\qquad
  \kappa_\wedge=\frac52\,\frac{nT_e}{m_e\Omega_e},
\]
with $\Omega_e$ the gyrofrequency. The parallel coefficient is the
unmagnetized one. The perpendicular coefficient is smaller by
$(\Omega\tau)^{-2}$, many orders of magnitude in a fusion plasma, because
the step of the random walk is the gyroradius and not the mean free path.
The third term is perpendicular to the gradient and does not depend on the
collision time. It is not dissipative. The stress tensor has the same
structure, with a non-dissipative \emph{gyroviscous} part.

For a closure, the consequences are two. The parallel direction carries
almost all the heat, so errors in the parallel closure and errors in the
numerical alignment with the field dominate. And the non-dissipative terms
are closures with a sign structure (they do no work) that a learned model
must preserve.

\subsection{Closing in phase space: drift kinetics and gyrokinetics}

The same small parameter can be used before taking moments. Averaging the
kinetic equation over the gyrophase removes one velocity coordinate and the
fast gyration time scale. In gyrokinetics the ordering is
\[
  \frac{\omega}{\Omega}\sim\frac{\rho}{L}\sim\frac{\delta f}{f}
  \sim\frac{k_\parallel}{k_\perp}\ll1,\qquad k_\perp\rho\sim1,
\]
and the result is a kinetic equation in five phase-space dimensions for
rings of charge \cite{BrizardHahm07}. This is a closure in the sense of
Section~\ref{sec:problem}: the retained state is the gyrophase-averaged
distribution, the discarded information is the gyrophase dependence, and
the limiting process is the ordering above. It is asymptotic, with a
conserved adiabatic invariant behind it. Gyrofluid models are moment
closures of gyrokinetics \cite{DorlandHammett93}.

\subsection{Fluid models as a ladder of closures}

Single-fluid models of a magnetized plasma differ in what they close and
how.
\begin{center}
\renewcommand{\arraystretch}{1.3}
\begin{tabularx}{\textwidth}{@{}P{3.3cm}P{5.0cm}Y@{}}
\toprule
Model & Closure of the pressure & Closure of the electric field \\
\midrule
Ideal MHD & scalar, adiabatic: $\dd(p/n^{5/3})/\dd t=0$
  & $\bE+\bu\times\bB=0$ \\
Resistive MHD & scalar, with heating
  & $\bE+\bu\times\bB=\eta\,\mathbf J$ \\
Hall and extended MHD & scalar or Braginskii
  & generalized Ohm's law \\
Double-adiabatic & $p_\perp$, $p_\parallel$, no heat flux & ideal \\
Landau fluid & $p_\perp$, $p_\parallel$, nonlocal parallel heat flux
  & ideal or generalized \\
\bottomrule
\end{tabularx}
\end{center}
The generalized Ohm's law is the electron momentum equation solved for the
electric field,
\begin{equation}
  \bE+\bu\times\bB=\eta\,\mathbf J+\frac{\mathbf J\times\bB}{ne}
   -\frac{\nabla\cdot\mathbf P_e}{ne}
   +\frac{m_e}{ne^2}\frac{\dd\mathbf J}{\dd t},
\end{equation}
and each MHD variant keeps a subset of the terms on the right. The
electron pressure tensor in it is unclosed. In collisionless reconnection
its off-diagonal part is what breaks the frozen-in condition, which is why
ten-moment models with a heat-flux closure are used there \cite{Wang15}.

\section{Strongly coupled and dusty plasmas}
\label{sec:strong}

The first application profile of the project is a dusty plasma: micron-size
grains, each carrying thousands of elementary charges, suspended in a weakly
ionized gas. The grains are few and individually tracked. Their coupling
parameter $\Gamma$ is large, so the weak-coupling closure of
Section~\ref{sec:plasma} does not apply to them. Three closures are
stacked.

\subsection{First closure: the effective force}

The full system contains electrons, ions, neutrals and fields. The
retained state is the grain positions and velocities. Eliminating the rest
leaves a force on each grain that the grains alone must determine.
\begin{itemize}
  \item The classical closure is an \emph{instantaneous pair law}: a
    screened Coulomb (Yukawa) potential
    $\phi(r)=Q^2e^{-r/\lambda_D}/(4\pi\varepsilon_0r)$, plus confinement
    and gas drag. The state of the grain system is then described by
    $\Gamma$ and the screening parameter $\kappa=a/\lambda_D$.
  \item What it discards: ion flow past a grain produces a wake, so the
    force between two grains is not reciprocal and not a gradient of a
    potential; grain charge depends on the local plasma and fluctuates; the
    wake takes time to form, which is memory; and the force on a grain
    depends on more than one neighbour.
\end{itemize}
In the terms of Section~\ref{sec:kinds} the candidates are an algebraic
pair closure, a many-body (graph) closure, and a memory closure, and the
project's decisive experiment nests them in that order. The eliminated
plasma has its own dynamics and is not in equilibrium, so the
fluctuation--dissipation relation \eqref{eq:gle} need not hold for the
effective grain dynamics. The data-driven force inference of the Burton
laboratory \cite{Yu25} is an instance of learning this closure.

\subsection{Second closure: correlations}

Given a pair law, the equilibrium hierarchy relates the pair distribution
$g(r)$ to the triplet distribution. Liquid-state theory calls its
approximations \emph{closures}, in exactly the sense used here
\cite{HansenMcDonald}.
\begin{itemize}
  \item \textbf{Kirkwood superposition:}
    $g_3(1,2,3)\approx g(1,2)\,g(2,3)\,g(1,3)$, the factorization of the
    triplet function into pairs.
  \item \textbf{Integral-equation closures.} The Ornstein--Zernike relation
    $h=c+n\,c*h$ defines the direct correlation function $c$ from
    $h=g-1$. It is an identity with two unknowns. A second relation closes
    it:
    \[
      \text{hypernetted chain:}\quad g=\exp(-\phi/T+h-c),
      \qquad
      \text{Percus--Yevick:}\quad c=(1-e^{\phi/T})\,g .
    \]
    The first is accurate for long-range, soft potentials such as Coulomb
    and Yukawa; the second for short-range, hard ones.
\end{itemize}
For the dynamics, the BBGKY closure $g_2\to0$ fails outright at
$\Gamma\gtrsim1$: the system is a liquid, and for the one-component plasma
it freezes at $\Gamma\approx175$. Kinetic descriptions that keep
correlations use them as input: local-field corrections to the mean field,
and the quasi-localized charge approximation, in which particles oscillate
in a slowly changing cage described by $g(r)$ \cite{GoldenKalman00}.

\subsection{Third closure: moments with memory}

At the fluid level a strongly coupled liquid supports shear waves at short
wavelength, which a Newtonian stress cannot. The standard closure is
viscoelastic: the stress relaxes over a time $\tau_m$,
\begin{equation}
  \Bigl(1+\tau_m\frac{\partial}{\partial t}\Bigr)\boldsymbol\pi
  =-\mu\bigl(\nabla\bu+\nabla\bu^{\mathsf T}
     -\tfrac23(\nabla\cdot\bu)\mathbf I\bigr),
\end{equation}
which is Cattaneo's law (Example~\ref{ex:cattaneo}) for momentum, a
one-pole memory kernel. It gives viscous behaviour for
$\omega\tau_m\ll1$ and elastic behaviour for $\omega\tau_m\gg1$. This
generalized-hydrodynamic model was applied to dusty plasmas by Kaw and Sen
\cite{KawSen98}. Here the memory is not a refinement: without it the
closure has the wrong qualitative behaviour.

\subsection{What transfers}

The dusty-plasma profile is where the discarded information is
\emph{observable}: grain trajectories give two- and three-particle
statistics directly, so the factorization defect $g_2$, and the size of the
memory and many-body parts of the force, can be measured instead of
assumed. That is the reason it is paired with the canonical tracks, in
which the same quantities are controlled by theorems.

\section{Radiation, neutrinos, and the supernova profile}
\label{sec:rad}

The second application profile, core-collapse supernovae, involves two
closures: one for neutrino transport, which is a moment closure with a long
history in radiation transport, and one for turbulence in a
spherically symmetric model, which is a Reynolds closure. Both are
recorded here because the first supplies the cleanest worked example of a
maximum-entropy closure and the second defines the benchmark target.

\subsection{The moment hierarchy of transport}

The specific intensity $I(\bx,\boldsymbol\Omega,t)$ obeys a linear kinetic
equation in which particles all move at speed $c$,
\[
  \frac1c\,\partial_tI+\boldsymbol\Omega\cdot\nabla I
  =-\sigma_tI+\frac{\sigma_s}{4\pi}\int I\,\dd\Omega'+S .
\]
Angular moments give the energy density $E=c^{-1}\int I\,\dd\Omega$, the
flux $\bF=\int\boldsymbol\Omega\,I\,\dd\Omega$ and the pressure tensor
$\mathbf P=c^{-1}\int\boldsymbol\Omega\boldsymbol\Omega\,I\,\dd\Omega$,
with
\begin{equation}
  \partial_tE+\nabla\cdot\bF=\text{sources},\qquad
  \frac1{c^2}\,\partial_t\bF+\nabla\cdot\mathbf P=-\frac{\sigma_t}{c}\,\bF .
\end{equation}
The pressure is unclosed. With $f=|\bF|/(cE)$ and
$\mathbf n=\bF/|\bF|$ the standard closures are written through an
\emph{Eddington factor} $\chi$,
\begin{equation}
  \mathbf P=\Bigl[\frac{1-\chi}{2}\,\mathbf I
     +\frac{3\chi-1}{2}\,\mathbf n\mathbf n\Bigr]E .
\end{equation}
Realizability requires $f\le1$ and $f^2\le\chi\le1$.

\subsection{The closures}

\begin{itemize}
  \item \textbf{Diffusion.} $\chi=\tfrac13$ and the flux equation in
    quasi-steady state: $\bF=-\tfrac{c}{3\sigma_t}\nabla E$. This is the
    Chapman--Enskog closure, valid when the mean free path $1/\sigma_t$ is
    short. In the opposite limit it gives $|\bF|>cE$: energy moving faster
    than light.
  \item \textbf{Flux-limited diffusion.} Keep the diffusion form with a
    limited coefficient,
    $\bF=-\tfrac{c\,\lambda(R)}{\sigma_t}\nabla E$ with
    $R=|\nabla E|/(\sigma_tE)$. The Levermore--Pomraning limiter
    \cite{LevermorePomraning81} is
    $\lambda(R)=(\coth R-1/R)/R$, which tends to $\tfrac13$ as $R\to0$ and
    to $1/R$ as $R\to\infty$, so that $|\bF|\to cE$. It is the transport
    analogue of the electron flux limiter of Section~\ref{sec:plasma}, with
    a derivation behind the limiter.
  \item \textbf{M1.} Evolve both $E$ and $\bF$ and close with a function
    $\chi(f)$.
  \item \textbf{Variable Eddington tensor.} Compute
    $\mathbf P/E$ from a separate, cheaper solution of the transport
    equation and use it to close the moment system.
\end{itemize}

\begin{example}[Two maximum-entropy Eddington factors]
\label{ex:m1}
Assume the intensity is isotropic in a frame moving with some velocity
$\beta c$, so that $I\propto(1-\beta\mu)^{-4}$ with $\mu$ the direction
cosine. This is how Levermore derived the closure \cite{Levermore84}, and
it coincides with maximizing the entropy of a photon gas subject to $E$
and $\bF$. Its moments give the closed form
\begin{equation}
  \chi(f)=\frac{3+4f^2}{5+2\sqrt{4-3f^2}} ,
  \label{eq:m1}
\end{equation}
with $\chi(0)=\tfrac13$ and $\chi(1)=1$. Maximizing the classical
(Boltzmann) entropy instead gives $I\propto e^{b\mu}$, with
$f=\coth b-1/b$ and $\chi=1-2f/b$ \cite{Minerbo78}. Both were checked by
quadrature: for example $\beta=0.4$ gives $f=0.5063$ and $\chi=0.4684$
from the integrals and from \eqref{eq:m1}.
\end{example}

The M1 closure is hyperbolic, realizable and exact in both limits, which
is everything Section~\ref{sec:moment} asked of a maximum-entropy closure.
Its known failure shows what the ansatz discards: for two beams crossing
head-on $f=0$ and the true $\chi$ is $1$, while \eqref{eq:m1} gives
$\tfrac13$. A function of $f$ alone cannot distinguish an isotropic field
from two opposed beams. Neutrino transport in supernova simulations has the
same structure with energy dependence and several species, and M1 is in
common use there.

\subsection{Turbulence in a one-dimensional supernova model}

A spherically symmetric model evolves angle averages. Let
$\avg{\cdot}_\Omega$ be the average over angles at fixed radius. Averaging
the three-dimensional equations and subtracting the one-dimensional
equations applied to the averaged state leaves
\begin{equation}
  C^\star_{\mathrm{SN}}[X]
  =\avg{F_{3\mathrm D}(X)}_\Omega-F_{1\mathrm D}\bigl(\avg{X}_\Omega\bigr),
\end{equation}
which is \eqref{eq:defect} for this projection. It decomposes into a
Reynolds stress (a turbulent pressure that helps support the shock), a
turbulent energy or enthalpy flux, production, dissipation and mixing. The
classical closure is mixing-length theory, the Prandtl model of
Section~\ref{sec:turb} with buoyancy as the driver, used in
one-dimensional explosion models \cite{CouchWarrenOConnor20}. What makes
this harder than the terrestrial case is that the flow is compressible,
passes through a moving shock, and never reaches a statistically steady
state, so history enters.

\section{Waves: wave kinetics and quasilinear theory}
\label{sec:waves}

\subsection{Wave turbulence}

For a weakly nonlinear dispersive system with four-wave interactions and
dispersion relation $\omega_k$, the random-phase closure gives the wave
kinetic equation for the spectrum $n_k$ \cite{ZLF92},
\begin{equation}
  \partial_tn_1\;\propto\;\lambda^2\!\int n_1n_2n_3n_4
  \Bigl(\frac1{n_1}+\frac1{n_2}-\frac1{n_3}-\frac1{n_4}\Bigr)
  \delta(k_1+k_2-k_3-k_4)\,\delta(\omega_1+\omega_2-\omega_3-\omega_4)
  \,\dd k_2\,\dd k_3\,\dd k_4 ,
\end{equation}
which has the structure of a Boltzmann collision integral: resonant
quartets of waves play the role of colliding pairs, and the kinetic time
scales as $\lambda^{-2}$. The five items are: hierarchy, the equations for
even-order correlations of Fourier amplitudes; limit, weak nonlinearity
with a large box, taken together; retained state, the spectrum; discarded,
phases and higher cumulants; defect, the difference between the measured
spectral transfer and the collision integral. Deng and Hani \cite{DH23}
proved the limit for the cubic nonlinear Schr\"odinger equation over the
full lifespan of the kinetic solution. It is the second canonical track of
the benchmark, chosen because the discarded information is spectral and
not particle-based.

\subsection{Quasilinear theory}

The plasma version closes the interaction between waves and the particle
distribution. Split $f=f_0+\delta f$, with $f_0$ the spatial average, in a
one-dimensional electrostatic plasma. The average of the Vlasov equation is
exact and unclosed,
\begin{equation}
  \partial_tf_0=-\frac{e}{m}\,\partial_v\avg{\delta E\,\delta f}.
\end{equation}
Quasilinear theory \cite{VVS62,DrummondPines62} closes it by using the
\emph{linear} response of $\delta f$ to $\delta E$ about the slowly
evolving $f_0$, which gives a diffusion equation in velocity,
\begin{equation}
  \partial_tf_0=\partial_v\bigl(D(v)\,\partial_vf_0\bigr),
  \qquad
  D(v)\propto\frac{e^2}{m^2}\int|E_k|^2\,\pi\,\delta(\omega_k-kv)\,\dd k,
  \qquad
  \partial_t|E_k|^2=2\gamma_k[f_0]\,|E_k|^2 ,
\end{equation}
with the growth rate $\gamma_k$ given by the slope of $f_0$ at the
resonant velocity. The constant of proportionality depends on the
normalization of the spectrum. The closure assumes many overlapping
resonances and a field that decorrelates before a particle is trapped.

Two connections to earlier sections. The delta function is the same
resonance as in Section~\ref{sec:collisionless}; quasilinear theory is its
nonlinear continuation for the mean distribution. And the natural repair
when the assumptions fail, resonance broadening \cite{Dupree66}, replaces
the delta function by a response function of finite width determined
self-consistently by the diffusion, which is a closure of the
direct-interaction type (Section~\ref{sec:statclosure}). Quasilinear
transport models with fitted saturation rules are in routine use for
turbulent transport in magnetized fusion plasmas, and they are closures in
the sense of Section~\ref{sec:turb}: models, calibrated against nonlinear
simulations.

\section{Learned closures}
\label{sec:learned}

Machine learning changes how a closure is represented and calibrated. It
does not remove the closure problem, and every structural question of the
previous sections applies unchanged. The literature chapter in
\texttt{shared/references/} treats the surrounding methods (PINNs, neural
operators, world models) at length; this section is restricted to closures
and to the design decisions that determine whether a learned closure can be
trusted. Reviews from the turbulence side are in
\cite{DuraisamyIaccarinoXiao19}.

\subsection{What is learned}

The first decision is the mathematical object. The options are the rows of
the table in Section~\ref{sec:kinds}, and the choice fixes what the model
can and cannot represent.
\begin{center}
\renewcommand{\arraystretch}{1.3}
\begin{tabularx}{\textwidth}{@{}P{3.6cm}P{5.2cm}Y@{}}
\toprule
Learned object & Form & Examples \\
\midrule
Coefficients of a classical closure
  & $\nu_t=g_\theta(\text{invariants})$
  & tensor-basis networks for the Reynolds stress \cite{Ling16} \\
The unclosed moment
  & $M_{N+1}=g_\theta(M_0,\dots,M_N)$
  & direct regression; no guarantee of hyperbolicity \\
Its gradient
  & $\partial_xM_{N+1}=\sum_ja_{j,\theta}\,\partial_xM_j$
  & hyperbolic moment closures \cite{Huang21,Christlieb25} \\
A nonlocal operator
  & $q=\mathcal W_\theta[T]$
  & learned Landau-fluid closures \cite{Ma20} \\
The choice of moments and the closure
  & encoder plus closed dynamics
  & learned generalized moments \cite{HanMaMaE19} \\
A memory kernel or auxiliary state
  & $K_\theta(s)$, or $\dot z=g_\theta(z,x)$
  & data-driven Mori--Zwanzig \cite{Lin21} \\
A collision operator
  & $C_\theta[f]$ in a conservative form
  & learned from molecular dynamics \cite{Zhao25} \\
A discretization-level correction
  & correction to a coarse solver step
  & learned corrections to coarse solvers \cite{Kochkov21} \\
\bottomrule
\end{tabularx}
\end{center}
Whole-solution surrogates (a network that maps initial data to the solution)
are not closures. They discard the reduced equation instead of completing
it, and in this project they are run as black-box controls.

\subsection{How structure enters}

A constraint can be imposed softly, as a penalty in the loss, or hard, by
construction. A soft constraint holds approximately on the training
distribution. A hard constraint holds everywhere and is the only kind that
supports a claim about behaviour in a long rollout.
\begin{itemize}
  \item \textbf{Conservation:} learn a flux, not a source, and apply a
    conservative difference to it.
  \item \textbf{Symmetry:} build inputs from invariants and outputs in a
    tensor basis (Galilean and rotational invariance).
  \item \textbf{Hyperbolicity:} parameterize the closure so that the flux
    Jacobian has real eigenvalues, for instance through a symmetrizer or
    through the eigenvalues themselves.
  \item \textbf{Dissipation:} parameterize a non-negative symbol or a
    positive semidefinite coefficient matrix.
  \item \textbf{Limits:} learn a correction to a classical closure that
    vanishes in the regime where the classical one is trusted,
    \begin{equation}
      q=q_{\mathrm{classical}}+\delta q_\theta ,
      \qquad \delta q_\theta\to0\ \text{as}\ k\lambda\to0 .
    \end{equation}
\end{itemize}
The last item is engineering principle 10 of this project: the classical
reduced operator stays explicit and the learned part supplies only an
identified correction. For the linear heat flux the results of
Section~\ref{sec:repo} make it concrete. In one velocity dimension the
correction satisfies $0\le\kF-\chi\le\kF Kk^2/(1+Kk^2)$, so its amplitude
and its long-wavelength scaling are known in advance. For a thermal drive
the exact symbol is $\int(m+k^2)^{-1}\dd\rho(m)$ with $\rho\ge0$, so a
finite sum $\sum_jw_j/(m_j+k^2)$ with non-negative weights has the right
analytic type, and two linear conditions on the weights fix the
conductivity and the leading deficit.

\subsection{How it is trained}

\begin{itemize}
  \item \textbf{A priori (offline):} regress the closure term on reference
    data. Cheap, and it optimizes the wrong quantity: the defect on
    reference states, not the behaviour of the closed model.
  \item \textbf{A posteriori (solver in the loop):} differentiate through
    the closed solver and penalize the error of the rolled-out solution
    \cite{ListChenThuerey22,Kochkov21}. It optimizes the right quantity,
    needs a differentiable solver, and couples the closure to one
    discretization.
  \item \textbf{Statistical targets:} fit long-time statistics when
    trajectories diverge chaotically and pointwise rollout error is
    meaningless.
\end{itemize}

\subsection{How it fails}

\begin{itemize}
  \item \textbf{The a priori to a posteriori gap.} A small defect on
    reference data, an unstable or biased closed model
    (Section~\ref{sec:turb}).
  \item \textbf{Distribution shift under rollout.} The closed model drifts
    to states outside the training set, where the closure is unconstrained.
  \item \textbf{Extrapolation in the regime parameter.} A closure trained
    at one collisionality, coupling strength or resolution is used at
    another. Figure~\ref{fig:between} and the electron results of
    Section~\ref{sec:repo} show how much a closure can differ at fixed
    conductivity.
  \item \textbf{Non-identifiability.} Different closures fit the same data
    when the data do not excite the directions in which they differ, such
    as frequencies near the collision rate.
  \item \textbf{Entanglement with numerics.} A correction trained at one
    grid and time step absorbs discretization error, and is then a
    different model at another resolution.
  \item \textbf{Leakage.} Test cases that influenced normalization,
    hyperparameters or early stopping.
\end{itemize}

\subsection{What must be declared}

A comparison between closures is meaningful only at a matched
\emph{information budget}: which equations, paired fields, ensembles,
cumulants, histories, closure targets and solver calls each method
received (see \texttt{model-bakeoff/docs/information\_budgets.md}). Access
to kinetic trajectories or to frequency responses is a different budget
from access to temperature profiles alone, and a method given more will
appear better for that reason.

\section{How a closure is judged}
\label{sec:judge}

\subsection{Structural properties}

These can be checked without reference data, and a closure that fails one
is wrong in a way that accuracy elsewhere does not repair.
\begin{center}
\renewcommand{\arraystretch}{1.3}
\begin{tabularx}{\textwidth}{@{}P{3.3cm}YP{4.4cm}@{}}
\toprule
Property & Statement & Where it failed above \\
\midrule
Conservation & mass, momentum and energy of the closed model are conserved
  & source-form corrections \\
Realizability & closed moments belong to a non-negative distribution; a
  stress is positive semidefinite; $|\bF|\le cE$
  & Boussinesq in strong strain; quasi-normal; diffusion at low opacity \\
Well-posedness & hyperbolic or parabolic, as the physics requires
  & Burnett; Grad far from equilibrium; unconstrained learned moments \\
Entropy & an entropy inequality, where one exists & truncated expansions \\
Invariance & Galilean, rotational, and correct tensor symmetry
  & networks acting on raw components \\
Dissipation sign & the closure removes energy where the exact one does, and
  non-dissipative terms do no work
  & negative conductivity at high $k$ \\
Causality & memory acts on the past only; the kernel is analytic in the
  upper half-plane
  & acausal fits in frequency \\
Reciprocity & Onsager symmetry of coupled fluxes & independently fitted
  thermal force and heat flux \\
\bottomrule
\end{tabularx}
\end{center}

\subsection{Limits}

A closure should reduce to the trusted result where the small parameter is
small, at the right rate. For the heat flux this means the conductivity
$\kF$, then the coefficient $K$ of the leading deficit
\eqref{eq:deficit}, then the collisionless constant. Each is a separate
test, and Section~\ref{sec:between} showed that passing the first and the
third does not imply the second.

\subsection{Accuracy, in two modes}

\begin{itemize}
  \item \textbf{A priori:} the closure term against the reference closure
    term on held-out reference states, reported pointwise, in integral
    norms, spectrally, and by regime.
  \item \textbf{A posteriori:} the closed model integrated in a solver,
    against the projected reference. Report conservation drift,
    realizability violations, stability, and long-rollout statistics, not
    only short-time pointwise error.
\end{itemize}
Inequality \eqref{eq:gronwall} is the reason both are required.

\subsection{Separating the sources of error}

A reported difference between a closed model and a reference contains
four contributions, and the benchmark requires them to be reported
separately (scientific principle 9).
\begin{enumerate}
  \item \textbf{Discretization error} of the reference and of the reduced
    solver. Measured by refinement.
  \item \textbf{Ensemble-estimation error} of any target that is a
    statistic. Measured by ensemble size and confidence intervals.
  \item \textbf{Asymptotic-limit error:} the difference between the
    reference at a finite value of the small parameter and the limiting
    model. Measured by sweeping the parameter.
  \item \textbf{Model-form error} of the closure itself.
\end{enumerate}
A learned correction can be credited only with a reduction of the fourth.
A failure to recover the expected limit is a solver, estimator or
specification problem until shown otherwise, and is not evidence for a
correction.

\subsection{A ladder of tests}

The heat-flux manuscript proposes four levels, each supporting a different
claim. A closure that passes one does not automatically pass the next.
\begin{enumerate}
  \item Fix the collision operator, the dimension, the moment convention and
    the drive, and compare $\chi(k)$ with exact formulas or proved
    intervals.
  \item Test the leading deficit coefficient as well as the limiting
    conductivity.
  \item Change the collision model at matched conductivity.
  \item Run an embedded initial-value or driven experiment, which tests the
    resolved evolution.
\end{enumerate}
For learned closures add a fifth: hold out the physics, not samples. An
operator-transfer test must assign the collision-rate family, the ion
charge, or the coupling regime to the test set before any normalization or
fitting.

\section{This project: what is proved here, and how the benchmark uses it}
\label{sec:repo}

\subsection{The heat-flux results}

The notes in \texttt{shared/docs/} and the manuscript
\texttt{plasma\_heat\_flux\_paper/manuscript.tex} study the function
$\chi(k,\omega)$ of Section~\ref{sec:between} for linear perturbations of a
Maxwellian, mostly at zero frequency. The manuscript is the current
statement; \texttt{summary\_of\_results.tex} predates several of its
sections. What follows is a reading guide, with the manuscript's
conditions. All of it is linear, for a homogeneous background, one Fourier
direction, and an unbounded or periodic domain. None of the proofs has been
checked by a second reader, and the Landau and electron results rest on
published coercivity estimates that are cited and not reproved.

\paragraph{Conventions.} The manuscript uses the energy flux, which in one
dimension is half the third central moment used in
Sections~\ref{sec:collisionless} and~\ref{sec:between}. Its
one-dimensional collisionless constant is $\sqrt{2/\pi}$, and its
conductivities are half those quoted above. Relative deficits are
unaffected (Appendix~\ref{app:conv}).

\paragraph{The response depends on the drive.} A response ratio is tied to
its input. Two drives are treated: a potential force (the Hammett--Perkins
setting) and an isotropic heating source. They share the conductivity and
differ at order $k^2$.

\begin{center}
\renewcommand{\arraystretch}{1.3}
\begin{tabularx}{\textwidth}{@{}P{3.7cm}YP{3.7cm}@{}}
\toprule
Setting & Established & Not established \\
\midrule
One velocity dimension, potential drive, any conserving symmetric
non-negative operator
  & dual variational principles;
    $\kF/(1+Kk^2)\le\chi(k)\le\kF$ at every $k$; $K$ is the exact leading
    deficit (under stated solvability hypotheses)
  & nonlinear or arbitrary-initial-data closure \\
BGK and Dougherty
  & closed forms \eqref{eq:theta-bgk}, \eqref{eq:theta-dougherty}; both
    limits; sharp leading deficits; exact frequency response for
    constant-$\nu$ BGK
  & equivalence to Coulomb collisions \\
Three dimensions, potential drive
  & a stress cross term appears; two-sided bounds at small $k$; leading
    coefficient; positivity when $k^2\kF\eta<3$
  & bounds at all $k$: they are false in general \\
Counterexamples
  & relaxation operators with a slow stress have $\chi>\kF$; with a slow
    heat flux as well, $\chi<0$
  & sharp thresholds \\
Like-particle Landau
  & weighted solvability; $|k|\chi\to3\sqrt{2/\pi}$ with rate
    $O(|k|^{-2/3}\log|k|)$; positivity at small and at large $k$
  & positivity at intermediate $k$; $\chi\le\kF$ at all $k$ \\
Three dimensions, thermal drive
  & $\kappa/(1+K_{\mathrm{th}}k^2)\le\chi_{\mathrm{th}}\le\kappa$ at every
    $k$, without momentum conservation; $\chi_{\mathrm{th}}$ is a Stieltjes
    function of $k^2$
  & the spectral measure beyond two moments \\
Electrons with ions of charge $Z$
  & the thermal theorems with the zero-current conductivity;
    $|k|\chi_{\mathrm{th}}\to\tfrac95\sqrt{2/\pi}$ for every $Z$; the
    normalized deficit grows linearly in $Z$ with slope $1.4305$
  & certified constants; a rate for the collisionless limit \\
Frequency dependence
  & the exact closure for data in the retained moments is a convolution in
    time with a positive-real kernel; bounds on the error of the
    quasi-static closure
  & the continuum statement for Landau \\
\bottomrule
\end{tabularx}
\end{center}

\paragraph{Numbers.} From polynomial Galerkin calculations: for
like-particle Landau collisions $\kF/\tau\simeq3.9502$ and
$K/\tau^2\simeq22.35$, a leading deficit about $7.45$ times that of the
Dougherty operator at matched conductivity. For electrons the normalized
deficit $K_{\mathrm{th}}/\kappa^2$ grows from about $4.3$ at $Z=1$ to $96$
at $Z=64$. In the embedded tests the exact quasi-static closure has
transient errors of up to $34\%$ in a one-dimensional model coupled to
sound waves, and stays within $2\%$ in the quasi-neutral decay of an
electron temperature perturbation at $Z=1$ and $|k|\lambda_e\le0.3$, where
Spitzer--H\"arm conduction is off by $42\%$.

\paragraph{What they are for.} These are analytic test targets for
classical and learned closures, with known answers at both ends and proved
intervals in between. They are not a general constitutive law.

\subsection{The benchmark tracks, in the five items}

\begin{center}
\renewcommand{\arraystretch}{1.3}
\small
\begin{tabularx}{\textwidth}{@{}P{2.3cm}P{2.6cm}P{2.5cm}P{2.5cm}Y@{}}
\toprule
Track & Hierarchy and limit & Retained & Discarded & Defect measured by \\
\midrule
A1 particles to Boltzmann
  & BBGKY; low density
  & one-particle distribution
  & cumulants, collision history
  & factorization error, recollision counts, Boltzmann residual \\
A2 waves to wave kinetics
  & correlation hierarchy; weak nonlinearity
  & spectrum
  & phases, spectral cumulants
  & spectral covariance error, kinetic residual \\
A3 Boltzmann to fluid
  & moment hierarchy; $\Kn\to0$
  & $n$, $\bu$, $T$
  & non-Maxwellian velocity structure
  & constitutive flux error, Hilbert remainder \\
B1 dusty plasma
  & three stacked closures (Section~\ref{sec:strong}); no small parameter
    assumed
  & grains; then $f_1$; then moments
  & plasma and wakes; correlations; velocity detail
  & force decomposition, factorization defect, flux residuals \\
B2 supernova
  & Reynolds decomposition of 3D onto 1D; none
  & radial profiles
  & nonradial motion, history
  & $C^\star_{\mathrm{SN}}$ by term \\
\bottomrule
\end{tabularx}
\end{center}
Track A is where the closure is a theorem and the asymptotic-limit error
can be separated from everything else. Track B is where the assumptions
fail and an augmented closure is needed. The heat-flux results sit between
them: a kinetic-to-moment closure with exact answers and a regime
parameter ($k\lambda$, and for electrons $Z$) that sweeps from the trusted
limit to its breakdown.

\subsection{The closure models compared}

The list in \texttt{mathematical\_closure\_program.md} maps onto the
sections of this tutorial as follows.
\begin{center}
\renewcommand{\arraystretch}{1.2}
\begin{tabularx}{\textwidth}{@{}YP{4.6cm}@{}}
\toprule
Model in the program & Section here \\
\midrule
Analytic and identifiable effective interaction laws & \ref{sec:strong} \\
Contextual or finite-memory graph particle models
  & \ref{sec:strong}, \ref{sec:memory} \\
Classical Markov kinetic or fluid closure
  & \ref{sec:asym}, \ref{sec:plasma}, \ref{sec:collisionless} \\
Explicit finite-cumulant or finite-moment closure
  & \ref{sec:hier}, \ref{sec:moment} \\
Finite-memory closure with kernels or auxiliary variables
  & \ref{sec:memory}, \ref{sec:between} \\
Structure-preserving learned correction & \ref{sec:learned} \\
Black-box surrogate as a control & \ref{sec:learned} \\
\bottomrule
\end{tabularx}
\end{center}

\section{Summary, exercises, and symbols}
\label{sec:summary}

\subsection{Side by side}

\begin{center}
\renewcommand{\arraystretch}{1.4}
\begin{tabularx}{\textwidth}{@{}P{2.7cm}P{2.8cm}P{2.8cm}Y@{}}
\toprule
  & Neutral gas & Fluid turbulence & Plasma \\
\midrule
  What relaxes the discarded variables
    & frequent short-range collisions
    & nothing fast; cascade
    & weak, speed-dependent collisions; phase mixing; gyromotion \\
  Small parameter
    & $\Kn$
    & none
    & several ($1/\Lambda$, $\lambda/L$, $\rho/L$), often not small \\
  Form of closure
    & local gradient law
    & modelled eddy viscosity
    & nonlocal in space, anisotropic, sometimes with memory \\
  Status
    & derived, with error control
    & modelled
    & derived only in limits; modelled or matched elsewhere \\
\bottomrule
\end{tabularx}
\end{center}

\begin{center}
\renewcommand{\arraystretch}{1.3}
\begin{tabularx}{\textwidth}{@{}P{3.0cm}P{3.6cm}YY@{}}
\toprule
Method & Needs & Gives & Typical failure \\
\midrule
Asymptotic (Chapman--Enskog) & a small parameter, a collision operator
  & derived coefficients, error estimate
  & no scale separation; higher orders ill posed \\
Ansatz (Grad, maximum entropy) & a choice of moments and shape
  & finite speeds, far-from-equilibrium range
  & realizability, hyperbolicity, cost \\
Gradient-diffusion model & a calibrated coefficient
  & robust, dissipative
  & wrong tensor structure, no backscatter, tuned \\
Memory & a kernel or auxiliary variables
  & exact in principle
  & kernel unknown; cost grows with memory length \\
Response matching (Hammett--Perkins) & the linear kinetic response
  & correct linear limits without collisions
  & linear only; irreversible; intermediate accuracy \\
Learned & reference data and a declared budget
  & flexible correction
  & extrapolation, instability when embedded \\
\bottomrule
\end{tabularx}
\end{center}

\paragraph{In one sentence.} In a gas the closure is justified because
collisions enforce local equilibrium below every scale that is resolved; in
a plasma that justification is often missing, so the discarded kinetic
information is not determined by the local fluid state, and the closure
becomes nonlocal, anisotropic, and tied to wave--particle resonance.

\subsection{Choosing a closure}

Four questions, in order.
\begin{enumerate}
  \item \emph{Is there a small parameter, and is it small where the model
    will be used?} If so, use the asymptotic closure and its error
    estimate, and test the limit.
  \item \emph{If not, what relaxes the discarded variables, and how fast
    compared with the resolved dynamics?} Fast and local: gradient law.
    Fast but long-range in space: nonlocal quasi-static closure.
    Comparable to the resolved time scale: retain the flux or a memory
    kernel.
  \item \emph{Which structural properties must hold exactly?} Build those
    in. Leave the rest to be tested.
  \item \emph{What reference data exist, and with what uncertainty?} That
    bounds what a learned correction can be credited with.
\end{enumerate}

\subsection{Exercises}

Short answers follow the list.

\begin{exercise}\label{ex:gron}
Derive \eqref{eq:gronwall}. Show that if $f$ satisfies the one-sided bound
$(x-y)\cdot(f(x)-f(y))\le-\gamma|x-y|^2$ with $\gamma>0$, the embedded error
is bounded by $\sup|d|/\gamma$ for all time.
\end{exercise}

\begin{exercise}
For \eqref{eq:fastslow}, expand the slow eigenvalue to $O(\varepsilon^2)$
and compare with the three closures. Which of the two first-order closures
has the smaller $O(\varepsilon^2)$ error, and does the answer depend on
$a$, $b$, $c$?
\end{exercise}

\begin{exercise}
Derive the moment hierarchy of Example~\ref{ex:langevin}. Verify
$m_4=\sigma^2/2$ in the stationary state without any closure, and both
closure values of $m_2$. Is the pair $(m_2,m_4)$ from the second closure
realizable?
\end{exercise}

\begin{exercise}
Verify \eqref{eq:ce-source}, and compute $\kappa$ and $\mu$ for BGK from
\eqref{eq:transport-forms}.
\end{exercise}

\begin{exercise}
For the temperature equation
$\partial_t\hat T=-k^2\chi(k)\hat T$ (equation \eqref{eq:T} with the flow
neglected), compare the growth rates for $\chi=\kF(1-Kk^2)$ and
$\chi=\kF/(1+Kk^2)$. What is the large-$k$ decay rate in the second case,
and how does it compare with the collisionless rate $\propto|k|$?
\end{exercise}

\begin{exercise}
Show that the eigenvalues of the truncated Jacobi matrix $J_N$ are the
roots of $\He_{N+1}$.
\end{exercise}

\begin{exercise}
Derive $T_k=2\zeta n_k/(\zeta+i\chi_1)$ and \eqref{eq:R3}, and obtain
$\chi_1=2/\sqrt\pi$.
\end{exercise}

\begin{exercise}
Prove the recurrence for $W_n$ and derive \eqref{eq:cN}. Show
$c_N\to1$.
\end{exercise}

\begin{exercise}
Show that the kernel $(2\ell)^{-1}e^{-|x|/\ell}$ has Fourier transform
$1/(1+k^2\ell^2)$, and verify \eqref{eq:lorentz}.
\end{exercise}

\begin{exercise}
Show that Cattaneo's law is equivalent to the exponential memory kernel of
Example~\ref{ex:cattaneo}, and find the dispersion relation of the
resulting temperature equation.
\end{exercise}

\begin{exercise}
Verify \eqref{eq:m1} for $I\propto(1-\beta\mu)^{-4}$ in slab geometry, and
show that it satisfies $f^2\le\chi\le1$.
\end{exercise}

\begin{exercise}\label{ex:smag}
Using the Kolmogorov spectrum $E(k)=C_K\varepsilon^{2/3}k^{-5/3}$ up to the
cutoff $\pi/\Delta$, and the approximation
$\avg{|\bar S|^3}\approx\avg{|\bar S|^2}^{3/2}$, derive
$C_s=\pi^{-1}(2/3C_K)^{3/4}$.
\end{exercise}

\begin{exercise}
Find the strain rate at which the Boussinesq closure \eqref{eq:boussinesq}
gives a negative normal stress in plane strain.
\end{exercise}

\begin{exercise}\label{ex:gamma}
Show $\Gamma=\tfrac13\Lambda^{-2/3}$.
\end{exercise}

\paragraph{Answers.}
\begin{enumerate}[label=\arabic*.]
  \item Multiply $\dot e=f(x_c)-f(\cP X)-d$ by $e$:
    $\tfrac12\tfrac{\dd}{\dd t}|e|^2\le-\gamma|e|^2+|d||e|$.
  \item $\lambda=(a+bc)-\varepsilon bc(a+bc)
    +\varepsilon^2bc(a+bc)(a+2bc)+\dots$ The polynomial closure has no
    $\varepsilon^2$ term; the rational one has
    $\varepsilon^2(bc)^2(a+bc)$. The errors are
    $\varepsilon^2|bc(a+bc)(a+2bc)|$ and $\varepsilon^2|bc(a+bc)^2|$, so the
    rational form is closer exactly when $a/(bc)>-\tfrac32$, as for $a=0$
    in the table.
  \item $\dot m_4=-4m_6+6\sigma^2m_2$ and $\dot m_2=-2m_4+\sigma^2$. Yes:
    $m_4=0.5\,\sigma^2\ge m_2^2=0.2\,\sigma^2$.
  \item $\kappa=\tfrac52nT/(m\nu)$, $\mu=nT/\nu$ (Example~\ref{ex:bgk}).
  \item The first grows like $k^4$ for $k^2>1/K$. The second decays at a
    rate tending to the constant $\kF/K$: bounded, where the collisionless
    rate grows like $|k|$. The rational form is well posed but
    underdamps short wavelengths.
  \item Expand $\det(v-J_N)$ along the last row; the determinants obey the
    three-term recurrence of $\He_n$.
  \item Section~\ref{sec:collisionless}.
  \item $N!!/(N-1)!!\sim\sqrt{\pi N/2}$ for even $N$ and
    $\sqrt{2N/\pi}$ for odd $N$ (Wallis).
  \item Direct integration; substitute $m=k^2t^2$ in \eqref{eq:lorentz}.
  \item $-i\omega(1-i\omega\tau)=-\kappa k^2$ in suitable units: two
    damped modes for $4\kappa k^2\tau<1$, propagating damped waves
    otherwise.
  \item The moments are elementary integrals of $(1-\beta\mu)^{-4}$;
    eliminate $\beta$. Checked numerically in \texttt{checks.py}.
  \item $\avg{|\bar S|^2}=2\int_0^{\pi/\Delta}k^2E\,\dd k
    =\tfrac32C_K\varepsilon^{2/3}(\pi/\Delta)^{4/3}$; set
    $\varepsilon=(C_s\Delta)^2\avg{|\bar S|^2}^{3/2}$.
  \item $\bar S_{11}>k/(3\nu_t)$.
  \item $e^2/(\varepsilon_0T)=1/(n\lambda_D^2)$ and $n=3/(4\pi a^3)$ give
    $\Gamma=a^2/(3\lambda_D^2)$, and $\Lambda=(\lambda_D/a)^3$.
\end{enumerate}

\subsection{Symbols}

\begin{center}
\renewcommand{\arraystretch}{1.15}
\begin{tabularx}{\textwidth}{@{}P{2.6cm}Y@{}}
\toprule
$X$, $x=\cP X$ & full and retained state \\
$d$, $e$ & a priori defect; a posteriori error \\
$f$, $f_M$, $h$ & distribution; local Maxwellian; relative perturbation,
   $f=M(1+h)$ \\
$n$, $\bu$, $T$, $p$ & density, flow, temperature (energy units), pressure \\
$\boldsymbol\pi$, $\bq$ & viscous stress, heat flux \\
$\vt$ & $\sqrt{T/m}$ (not $\sqrt{2T/m}$) \\
$\Lambda$ (operator) & linearized collision operator, non-negative \\
$\Lambda$ (number) & particles per Debye sphere; $\ln\Lambda$ the Coulomb
   logarithm \\
$\nu$, $\tau$, $\lambda$ & collision frequency, collision time, mean free
   path \\
$\Kn$, $\Gamma$, $\rho$ & Knudsen number, coupling parameter, gyroradius \\
$\kappa$, $\kF$, $\mu$ & conductivity; its Chapman--Enskog value;
   viscosity \\
$\chi(k,\omega)$ & heat-flux multiplier, $\hat q=-ik\chi\hat T$ \\
$K$ & leading deficit, $\chi=\kF(1-Kk^2+\dots)$ \\
$\theta$, $y$, $s$ & $|k|\chi$; $\nu/|k|$; $|k|\kF/(n_0\vt)$ \\
$\zeta$, $Z$, $R$ & $\omega/(\sqrt2|k|\vt)$; plasma dispersion function;
   $1+\zeta Z$ \\
$a_n$, $\psi_n$, $\He_n$ & Hermite coefficients, orthonormal Hermite
   functions, Hermite polynomials \\
$\Hil$ & Hilbert transform, multiplier $-i\,\sgn(k)$ \\
$R_{ij}$, $\tau_{ij}$, $\nu_t$ & Reynolds stress, subgrid stress, eddy
   viscosity \\
$E$, $\bF$, $\mathbf P$, $\chi(f)$ & radiation energy density, flux,
   pressure; Eddington factor (Section~\ref{sec:rad} only) \\
\bottomrule
\end{tabularx}
\end{center}


\appendix
\section{Conventions and conversions}
\label{app:conv}

\begin{center}
\renewcommand{\arraystretch}{1.3}
\begin{tabularx}{\textwidth}{@{}P{3.6cm}YY@{}}
\toprule
Quantity & This tutorial and the Hammett--Perkins papers
  & Elsewhere \\
\midrule
Thermal speed & $\vt=\sqrt{T_0/m}$
  & $v_{\mathrm{th}}=\sqrt{2T_0/m}$ in \cite{Hunana19}: coefficients of
    $\vt$ are divided by $\sqrt2$ per power \\
Heat flux, one velocity dimension & third central moment,
  $q=\avg{\He_3h}$
  & energy flux $\tfrac12\avg{\He_3h}$ in the heat-flux manuscript:
    $q$, $\chi$, $\kF$, $\theta$ are halved; $K$ is unchanged; $K/\kF^2$ is
    multiplied by four \\
Collisionless constant, one dimension & $\sqrt{8/\pi}$
  & $\sqrt{2/\pi}$ in the manuscript \\
Collisionless constant, three dimensions & not derived here
  & $3\sqrt{2/\pi}$ (potential drive) and $\tfrac95\sqrt{2/\pi}$ (thermal
    drive) in the manuscript \\
Fourier modes & $e^{ikz-i\omega t}$ & \\
Hilbert transform & multiplier $-i\,\sgn(k)$
  & opposite sign in several Landau-fluid papers \\
Temperature & energy units, $k_B=1$ & \\
\bottomrule
\end{tabularx}
\end{center}
Comparisons across velocity dimensions made from mixed conventions are not
meaningful; the manuscript makes physical comparisons only within a fixed
dimension, and so does this tutorial.

\section{Hermite relations used}
\label{app:hermite}

Probabilists' Hermite polynomials: $\He_0=1$, $\He_1=v$,
$\He_{n+1}=v\,\He_n-n\,\He_{n-1}$, with $\avg{\He_m\He_n}=n!\,\delta_{mn}$
for the unit Gaussian and $\He_n'=n\,\He_{n-1}$. The orthonormal functions
$\psi_n=\He_n/\sqrt{n!}$ satisfy
$v\,\psi_n=\sqrt{n+1}\,\psi_{n+1}+\sqrt n\,\psi_{n-1}$. The operator
$-\partial_v^2+v\,\partial_v$ has eigenfunctions $\He_n$ with eigenvalues
$n$, which is why the Lenard--Bernstein and Dougherty operators are
diagonal in this basis and damp high $n$ strongly. Values at the origin:
$\He_{2j}(0)=(-1)^j(2j-1)!!$.

\section{What was checked numerically}
\label{app:checks}

Run recorded in \texttt{\rundir/checks.json}; command
\texttt{python3 checks.py --output \rundir}. No random numbers are used.
\begin{center}
\renewcommand{\arraystretch}{1.25}
\begin{tabularx}{\textwidth}{@{}P{4.2cm}YP{3.2cm}@{}}
\toprule
Claim & Compared with & Agreement \\
\midrule
Fast--slow table & eigenvalues of the $2\times2$ matrix & exact \\
Langevin moments & quadrature of the stationary density & $10^{-15}$ \\
BGK $\kappa\nu=\tfrac52$, $\mu\nu=1$, 1D $\kF\nu=3$
  & exact Gaussian moments in rational arithmetic & exact \\
Truncation nodes, $N=8$ & Gauss--Hermite nodes from NumPy
  & $3\times10^{-15}$ \\
Eddington factors & quadrature of the ansatz & $10^{-15}$ \\
Hammett--Perkins: orders of matching
  & $R$ from the Faddeeva function: error $\propto\zeta^2$ ($R_3$),
    $\zeta^3$ ($R_4$) as $\zeta\to0$; $\zeta^{-5}$, $\zeta^{-6}$ as
    $\zeta\to\infty$
  & orders $2.00$, $3.00$, $5.01$, $6.02$ \\
$\sup|R_n-R|$ & grid of 24001 points on $0\le\zeta\le12$ & $0.344$, $0.059$ \\
$c_N$, $N=2,\dots,8$ & quadrature of $W_n$ at $z=10^{-4}i$
  & $2\times10^{-4}$ \\
$\theta(y)$ for BGK & table in
  \texttt{plasma\_closure\_gaps\_and\_bgk\_theorem.tex}; Hermite--Galerkin
  solve, 1200 modes
  & 4 digits; $10^{-14}$ \\
$\theta_1(y)$ for Dougherty & Hermite--Galerkin solve, 1200 modes
  & $10^{-15}$ \\
Exact over interpolant & 81 points, $10^{-2}\le s\le10^2$
  & ranges quoted in Section~\ref{sec:between} \\
\bottomrule
\end{tabularx}
\end{center}
Not checked numerically: any statement about three velocity dimensions,
about the Landau operator, or about a speed-dependent collision frequency.
Those are quoted from the notes named in the text.


\paragraph{Note on the references.} Entries were checked on 24 November 2026 against publisher, DOI or
catalogue records for authors, title, journal, volume, first page and year,
or were taken from the audited bibliographies of the heat-flux manuscript
and the literature chapter in this repository. Three were confirmed only in
part: \cite{Bobylev82} through a secondary reference list,
\cite{VVS62} without a primary record of its title, and
\cite{Lanford75} without a record of the volume number. The check covers
the citation data. It does not cover whether each source says what the
text attributes to it.

\begin{thebibliography}{99}

% --- kinetic theory and hydrodynamic limits ---
\bibitem{ChapmanCowling}
S.~Chapman and T.~G. Cowling,
\emph{The Mathematical Theory of Non-uniform Gases},
3rd ed., Cambridge University Press (1970).
\bibitem{BGK54}
P.~L. Bhatnagar, E.~P. Gross, and M.~Krook,
\emph{A model for collision processes in gases. I. Small amplitude processes
in charged and neutral one-component systems},
Phys.\ Rev.\ \textbf{94}, 511 (1954).
\bibitem{Grad49}
H.~Grad,
\emph{On the kinetic theory of rarefied gases},
Commun.\ Pure Appl.\ Math.\ \textbf{2}, 331 (1949).
\bibitem{Levermore96}
C.~D. Levermore,
\emph{Moment closure hierarchies for kinetic theories},
J.\ Stat.\ Phys.\ \textbf{83}, 1021 (1996).
\bibitem{Junk98}
M.~Junk,
\emph{Domain of definition of Levermore's five-moment system},
J.\ Stat.\ Phys.\ \textbf{93}, 1143 (1998).
\bibitem{StruchtrupTorrilhon03}
H.~Struchtrup and M.~Torrilhon,
\emph{Regularization of Grad's 13 moment equations: derivation and linear
analysis},
Phys.\ Fluids \textbf{15}, 2668 (2003).
\bibitem{Struchtrup05}
H.~Struchtrup,
\emph{Macroscopic Transport Equations for Rarefied Gas Flows},
Springer (2005).
\bibitem{McGraw97}
R.~McGraw,
\emph{Description of aerosol dynamics by the quadrature method of moments},
Aerosol Sci.\ Technol.\ \textbf{27}, 255 (1997).
\bibitem{Bobylev82}
A.~V. Bobylev,
\emph{The Chapman--Enskog and Grad methods for solving the Boltzmann
equation},
Sov.\ Phys.\ Dokl.\ \textbf{27}, 29 (1982).
\bibitem{GorbanKarlin14}
A.~N. Gorban and I.~Karlin,
\emph{Hilbert's 6th problem: exact and approximate hydrodynamic manifolds
for kinetic equations},
Bull.\ Amer.\ Math.\ Soc.\ \textbf{51}, 187 (2014).
\bibitem{KogelbauerKarlin24}
F.~Kogelbauer and I.~Karlin,
\emph{Exact hydrodynamic manifolds for the linear Boltzmann BGK equation I:
spectral theory},
Contin.\ Mech.\ Thermodyn.\ \textbf{36}, 1685 (2024).
\bibitem{Lanford75}
O.~E. Lanford,
\emph{Time evolution of large classical systems},
in Dynamical Systems, Theory and Applications,
Lecture Notes in Physics \textbf{38}, Springer (1975), p.~1.
\bibitem{Caflisch80}
R.~E. Caflisch,
\emph{The fluid dynamic limit of the nonlinear Boltzmann equation},
Commun.\ Pure Appl.\ Math.\ \textbf{33}, 651 (1980).
\bibitem{BGL91}
C.~Bardos, F.~Golse, and D.~Levermore,
\emph{Fluid dynamic limits of kinetic equations. I. Formal derivations},
J.\ Stat.\ Phys.\ \textbf{63}, 323 (1991).
\bibitem{GSR04}
F.~Golse and L.~Saint-Raymond,
\emph{The Navier--Stokes limit of the Boltzmann equation for bounded
collision kernels},
Invent.\ Math.\ \textbf{155}, 81 (2004).
\bibitem{DHM24}
Y.~Deng, Z.~Hani, and X.~Ma,
\emph{Long time derivation of the Boltzmann equation from hard sphere
dynamics},
\href{https://arxiv.org/abs/2408.07818}{arXiv:2408.07818}.
\bibitem{DHM25}
Y.~Deng, Z.~Hani, and X.~Ma,
\emph{Hilbert's sixth problem: derivation of fluid equations via
Boltzmann's kinetic theory},
\href{https://arxiv.org/abs/2503.01800}{arXiv:2503.01800}.
\bibitem{DH23}
Y.~Deng and Z.~Hani,
\emph{Long time justification of wave turbulence theory},
\href{https://arxiv.org/abs/2311.10082}{arXiv:2311.10082}.

% --- turbulence ---
\bibitem{Pope}
S.~B. Pope,
\emph{Turbulent Flows},
Cambridge University Press (2000).
\bibitem{LaunderSpalding74}
B.~E. Launder and D.~B. Spalding,
\emph{The numerical computation of turbulent flows},
Comput.\ Methods Appl.\ Mech.\ Eng.\ \textbf{3}, 269 (1974).
\bibitem{Smagorinsky63}
J.~Smagorinsky,
\emph{General circulation experiments with the primitive equations. I. The
basic experiment},
Mon.\ Weather Rev.\ \textbf{91}, 99 (1963).
\bibitem{Germano91}
M.~Germano, U.~Piomelli, P.~Moin, and W.~H. Cabot,
\emph{A dynamic subgrid-scale eddy viscosity model},
Phys.\ Fluids A \textbf{3}, 1760 (1991).
\bibitem{Lilly92}
D.~K. Lilly,
\emph{A proposed modification of the Germano subgrid-scale closure method},
Phys.\ Fluids A \textbf{4}, 633 (1992).
\bibitem{Clark79}
R.~A. Clark, J.~H. Ferziger, and W.~C. Reynolds,
\emph{Evaluation of subgrid-scale models using an accurately simulated
turbulent flow},
J.\ Fluid Mech.\ \textbf{91}, 1 (1979).
\bibitem{Orszag70}
S.~A. Orszag,
\emph{Analytical theories of turbulence},
J.\ Fluid Mech.\ \textbf{41}, 363 (1970).
\bibitem{Kraichnan59}
R.~H. Kraichnan,
\emph{The structure of isotropic turbulence at very high Reynolds numbers},
J.\ Fluid Mech.\ \textbf{5}, 497 (1959).

% --- memory ---
\bibitem{Mori65}
H.~Mori,
\emph{Transport, collective motion, and Brownian motion},
Prog.\ Theor.\ Phys.\ \textbf{33}, 423 (1965).
\bibitem{Zwanzig61}
R.~Zwanzig,
\emph{Memory effects in irreversible thermodynamics},
Phys.\ Rev.\ \textbf{124}, 983 (1961).
\bibitem{Zwanzig}
R.~Zwanzig,
\emph{Nonequilibrium Statistical Mechanics},
Oxford University Press (2001).
\bibitem{ChorinHaldKupferman02}
A.~J. Chorin, O.~H. Hald, and R.~Kupferman,
\emph{Optimal prediction with memory},
Physica D \textbf{166}, 239 (2002).
\bibitem{Lin21}
K.~K. Lin and F.~Lu,
\emph{Data-driven model reduction, Wiener projections, and the
Koopman--Mori--Zwanzig formalism},
J.\ Comput.\ Phys.\ \textbf{424}, 109864 (2021).

% --- plasma: collisional ---
\bibitem{HelanderSigmar02}
P.~Helander and D.~J. Sigmar,
\emph{Collisional Transport in Magnetized Plasmas},
Cambridge University Press (2002).
\bibitem{Landau36}
L.~D. Landau,
\emph{Die kinetische Gleichung f\"ur den Fall Coulombscher Wechselwirkung},
Phys.\ Z.\ Sowjetunion \textbf{10}, 154 (1936).
\bibitem{Lenard60}
A.~Lenard,
\emph{On Bogoliubov's kinetic equation for a spatially homogeneous plasma},
Ann.\ Phys.\ \textbf{10}, 390 (1960).
\bibitem{Balescu60}
R.~Balescu,
\emph{Irreversible processes in ionized gases},
Phys.\ Fluids \textbf{3}, 52 (1960).
\bibitem{LenardBernstein58}
A.~Lenard and I.~B. Bernstein,
\emph{Plasma oscillations with diffusion in velocity space},
Phys.\ Rev.\ \textbf{112}, 1456 (1958).
\bibitem{Dougherty64}
J.~P. Dougherty,
\emph{Model Fokker--Planck equation for a plasma and its solution},
Phys.\ Fluids \textbf{7}, 1788 (1964).
\bibitem{SH53}
L.~Spitzer and R.~H\"arm,
\emph{Transport phenomena in a completely ionized gas},
Phys.\ Rev.\ \textbf{89}, 977 (1953).
\bibitem{Braginskii65}
S.~I. Braginskii,
\emph{Transport processes in a plasma},
in Reviews of Plasma Physics, Vol.~1, p.~205 (1965).
\bibitem{LMV83}
J.~F. Luciani, P.~Mora, and J.~Virmont,
\emph{Nonlocal heat transport due to steep temperature gradients},
Phys.\ Rev.\ Lett.\ \textbf{51}, 1664 (1983).
\bibitem{ES91}
E.~M. Epperlein and R.~W. Short,
\emph{A practical nonlocal model for electron heat transport in laser
plasmas},
Phys.\ Fluids B \textbf{3}, 3092 (1991).
\bibitem{SNB00}
G.~P. Schurtz, Ph.~D. Nicola\"i, and M.~Busquet,
\emph{A nonlocal electron conduction model for multidimensional radiation
hydrodynamics codes},
Phys.\ Plasmas \textbf{7}, 4238 (2000).
\bibitem{JiHeld14}
J.-Y. Ji and E.~D. Held,
\emph{Electron parallel closures for arbitrary collisionality},
Phys.\ Plasmas \textbf{21}, 122116 (2014).
\bibitem{Brodrick17}
J.~P. Brodrick et al.,
\emph{Testing nonlocal models of electron thermal conduction for magnetic
and inertial confinement fusion applications},
Phys.\ Plasmas \textbf{24}, 092309 (2017).

% --- plasma: collisionless and magnetized ---
\bibitem{Landau46}
L.~D. Landau,
\emph{On the vibrations of the electronic plasma},
J.\ Phys.\ USSR \textbf{10}, 25 (1946).
\bibitem{CGL56}
G.~F. Chew, M.~L. Goldberger, and F.~E. Low,
\emph{The Boltzmann equation and the one-fluid hydromagnetic equations in
the absence of particle collisions},
Proc.\ R.\ Soc.\ Lond.\ A \textbf{236}, 112 (1956).
\bibitem{Gould67}
R.~W. Gould, T.~M. O'Neil, and J.~H. Malmberg,
\emph{Plasma wave echo},
Phys.\ Rev.\ Lett.\ \textbf{19}, 219 (1967).
\bibitem{HP90}
G.~W. Hammett and F.~W. Perkins,
\emph{Fluid moment models for Landau damping with application to the
ion-temperature-gradient instability},
Phys.\ Rev.\ Lett.\ \textbf{64}, 3019 (1990).
\bibitem{HDP92}
G.~W. Hammett, W.~Dorland, and F.~W. Perkins,
\emph{Fluid models of phase mixing, Landau damping, and nonlinear
gyrokinetic dynamics},
Phys.\ Fluids B \textbf{4}, 2052 (1992).
\bibitem{DorlandHammett93}
W.~Dorland and G.~W. Hammett,
\emph{Gyrofluid turbulence models with kinetic effects},
Phys.\ Fluids B \textbf{5}, 812 (1993).
\bibitem{Snyder97}
P.~B. Snyder, G.~W. Hammett, and W.~Dorland,
\emph{Landau fluid models of collisionless magnetohydrodynamics},
Phys.\ Plasmas \textbf{4}, 3974 (1997).
\bibitem{BrizardHahm07}
A.~J. Brizard and T.~S. Hahm,
\emph{Foundations of nonlinear gyrokinetic theory},
Rev.\ Mod.\ Phys.\ \textbf{79}, 421 (2007).
\bibitem{DJU14}
A.~M. Dimits, I.~Joseph, and M.~V. Umansky,
\emph{A fast non-Fourier method for Landau-fluid operators},
Phys.\ Plasmas \textbf{21}, 055907 (2014).
\bibitem{Wang15}
L.~Wang, A.~H. Hakim, A.~Bhattacharjee, and K.~Germaschewski,
\emph{Comparison of multi-fluid moment models with particle-in-cell
simulations of collisionless magnetic reconnection},
Phys.\ Plasmas \textbf{22}, 012108 (2015).
\bibitem{Hunana19}
P.~Hunana et al.,
\emph{An introductory guide to fluid models with anisotropic temperatures.
Part 2. Kinetic theory, Pad\'e approximants and Landau fluid closures},
J.\ Plasma Phys.\ \textbf{85}, 205850603 (2019);
\href{https://arxiv.org/abs/1901.09360}{arXiv:1901.09360}.

% --- strong coupling ---
\bibitem{HansenMcDonald}
J.-P. Hansen and I.~R. McDonald,
\emph{Theory of Simple Liquids},
4th ed., Academic Press (2013).
\bibitem{GoldenKalman00}
K.~I. Golden and G.~J. Kalman,
\emph{Quasilocalized charge approximation in strongly coupled plasma
physics},
Phys.\ Plasmas \textbf{7}, 14 (2000).
\bibitem{KawSen98}
P.~K. Kaw and A.~Sen,
\emph{Low frequency modes in strongly coupled dusty plasmas},
Phys.\ Plasmas \textbf{5}, 3552 (1998).
\bibitem{Yu25}
W.~Yu, E.~Abdelaleem, I.~Nemenman, and J.~C. Burton,
\emph{Physics-tailored machine learning reveals unexpected physics in dusty
plasmas},
Proc.\ Natl.\ Acad.\ Sci.\ \textbf{122}, e2505725122 (2025).

% --- radiation and supernovae ---
\bibitem{LevermorePomraning81}
C.~D. Levermore and G.~C. Pomraning,
\emph{A flux-limited diffusion theory},
Astrophys.\ J.\ \textbf{248}, 321 (1981).
\bibitem{Levermore84}
C.~D. Levermore,
\emph{Relating Eddington factors to flux limiters},
J.\ Quant.\ Spectrosc.\ Radiat.\ Transfer \textbf{31}, 149 (1984).
\bibitem{Minerbo78}
G.~N. Minerbo,
\emph{Maximum entropy Eddington factors},
J.\ Quant.\ Spectrosc.\ Radiat.\ Transfer \textbf{20}, 541 (1978).
\bibitem{CouchWarrenOConnor20}
S.~M. Couch, M.~L. Warren, and E.~P. O'Connor,
\emph{Simulating turbulence-aided neutrino-driven core-collapse supernova
explosions in one dimension},
Astrophys.\ J.\ \textbf{890}, 127 (2020).

% --- waves ---
\bibitem{ZLF92}
V.~E. Zakharov, V.~S. L'vov, and G.~Falkovich,
\emph{Kolmogorov Spectra of Turbulence I: Wave Turbulence},
Springer (1992).
\bibitem{VVS62}
A.~A. Vedenov, E.~P. Velikhov, and R.~Z. Sagdeev,
\emph{Quasi-linear theory of plasma oscillations},
Nucl.\ Fusion Suppl.\ \textbf{2}, 465 (1962).
\bibitem{DrummondPines62}
W.~E. Drummond and D.~Pines,
\emph{Non-linear stability of plasma oscillations},
Nucl.\ Fusion Suppl.\ \textbf{3}, 1049 (1962).
\bibitem{Dupree66}
T.~H. Dupree,
\emph{A perturbation theory for strong plasma turbulence},
Phys.\ Fluids \textbf{9}, 1773 (1966).

% --- learned closures ---
\bibitem{DuraisamyIaccarinoXiao19}
K.~Duraisamy, G.~Iaccarino, and H.~Xiao,
\emph{Turbulence modeling in the age of data},
Annu.\ Rev.\ Fluid Mech.\ \textbf{51}, 357 (2019).
\bibitem{Ling16}
J.~Ling, A.~Kurzawski, and J.~Templeton,
\emph{Reynolds averaged turbulence modelling using deep neural networks
with embedded invariance},
J.\ Fluid Mech.\ \textbf{807}, 155 (2016).
\bibitem{HanMaMaE19}
J.~Han, C.~Ma, Z.~Ma, and W.~E,
\emph{Uniformly accurate machine learning-based hydrodynamic models for
kinetic equations},
Proc.\ Natl.\ Acad.\ Sci.\ \textbf{116}, 21983 (2019).
\bibitem{Ma20}
C.~Ma, B.~Zhu, X.-Q. Xu, and W.~Wang,
\emph{Machine learning surrogate models for Landau fluid closure},
Phys.\ Plasmas \textbf{27}, 042502 (2020).
\bibitem{Huang21}
J.~Huang, Y.~Cheng, A.~J. Christlieb, and L.~F. Roberts,
\emph{Machine learning moment closure models for the radiative transfer
equation I: directly learning a gradient based closure},
\href{https://arxiv.org/abs/2105.05690}{arXiv:2105.05690}.
\bibitem{Christlieb25}
A.~J. Christlieb, M.~Ding, J.~Huang, and N.~A. Krupansky,
\emph{Hyperbolic machine learning moment closures for the BGK equations},
Multiscale Model.\ Simul.\ \textbf{23}, 187 (2025).
\bibitem{Zhao25}
Y.~Zhao, J.~W. Burby, A.~Christlieb, and H.~Lei,
\emph{Data-driven construction of a generalized kinetic collision operator
from molecular dynamics},
\href{https://arxiv.org/abs/2503.24208}{arXiv:2503.24208}.
\bibitem{Kochkov21}
D.~Kochkov, J.~A. Smith, A.~Alieva, Q.~Wang, M.~P. Brenner, and S.~Hoyer,
\emph{Machine learning--accelerated computational fluid dynamics},
Proc.\ Natl.\ Acad.\ Sci.\ \textbf{118}, e2101784118 (2021).
\bibitem{ListChenThuerey22}
B.~List, L.-W. Chen, and N.~Thuerey,
\emph{Learned turbulence modelling with differentiable fluid solvers:
physics-based loss functions and optimisation horizons},
J.\ Fluid Mech.\ \textbf{949}, A25 (2022).

\end{thebibliography}


\end{document}
